\documentclass[11pt]{article}
\usepackage[letterpaper,margin=1in]{geometry}
\usepackage[T1]{fontenc}
\usepackage{lmodern,microtype}
\usepackage{amsmath,amssymb,amsthm,mathtools}
\usepackage{booktabs,longtable,array}
\usepackage{xcolor,xurl}
\usepackage[colorlinks=true,linkcolor=blue!45!black,citecolor=blue!45!black,urlcolor=blue!45!black]{hyperref}
\hypersetup{pdftitle={Quantifier elimination, variable orders, and program proofs},pdfauthor={Charlie Cruz}}
\setlength{\emergencystretch}{2em}
\newtheorem{theorem}{Theorem}[section]
\newtheorem{lemma}[theorem]{Lemma}
\newtheorem{proposition}[theorem]{Proposition}
\newtheorem{corollary}[theorem]{Corollary}
\theoremstyle{definition}
\newtheorem{definition}[theorem]{Definition}
\theoremstyle{remark}
\newtheorem{remark}[theorem]{Remark}
\title{\textbf{Quantifier elimination, variable orders,\protect\\ and program proofs}\\[0.45em]
\large Further logical consequences of integral counting}
\author{Charlie Cruz}
\date{19 November 2026}
\begin{document}
\maketitle
\begin{abstract}
The integral contraction recursion of~\cite{Cruz26} has more structure
than its complexity bound records. Every merge in the recursion makes
one global row/column choice. Grouping the expanded formula by these
choices writes it as a sum of at most $h$ ordered formulas of total
size $2^{O(s/\sqrt h)}$, for every power of two $h$ with $2^{3h}\le s$.
The $h$ variable orders can be prescribed independently of the network.
Combined with disjoint counting sparsification, this gives
subexponential weighted-model formulas and integral multilinear IPS
refutations that use any number of orders tending to infinity.

A 2026 construction of fast codes and a specialization of Homeister's
lower-bound argument show that this is qualitatively tight for an
explicit family of unambiguous existential 3-CNFs: eliminating the
witnesses into an XOR of ordered parity programs has subexponential
minimum size if and only if the permitted number of orders tends to
infinity. Separately, the time--space simulation of~\cite{Cruz26}
implies a lower bound for equivalence of ordinary binary regular
expressions: no nondeterministic multitape algorithm decides it in
time $o(n\sqrt{\log n})$, where $n$ is the full syntax length. As a
consequence, fully written tree-like Kleene-algebra proofs cannot have
uniformly near-linear total symbol size. We also prove a restricted
same-arity separation in descriptive complexity, and we isolate the
unproved reflection step that would yield subexponential Frege
refutations.
\end{abstract}

\tableofcontents

\section{Introduction}
\label{sec:premises}

Cruz~\cite{Cruz26} simulates nondeterministic multitape time $t$ in
deterministic space $O(t/\sqrt{\log t})$ by counting computations with
an integral arithmetic formula. The formula comes from a contraction
recursion on a sparse network. The network is first grouped into an
initial clique of $q=2^h$ nodes, and each subsequent merge contracts
the clique along its rows or along its columns. The preceding
note~\cite{PreviousNote} lifts this construction to symbolic weights
with unique owners and derives first-failed-clause IPS certificates
from it.

This paper starts from one structural observation. Each merge makes
a single \emph{global} row/column choice, and all these choices are
made before any table entry is unfolded. Grouping the expanded formula
by its history of cuts therefore writes it as a sum of at most $h$
formulas, each of which respects one variable order. The orders are
determined by the cut histories alone, and with a suitable placement
of the weight owners they do not depend on the network. Most results
below are consequences of this observation; the remaining ones use the
simulation theorem of~\cite{Cruz26} directly.

\subsection{Main results}

\paragraph{Result 1: a tunable, universal order list.}
For a Boolean circuit whose gates and input vertices total $s$, and
every integer order budget $k\ge1$, the compiler gives at most $k$
ordered parity programs of total size
\[
  2^{O\!\left(s/\sqrt{\min\{k,\log(s+2)\}}\right)}.
\]
The hidden constant does not depend on $k$. The permitted list of
orders is fixed from the input-universe size and the recursion
parameter, independently of the circuit. This is a nonuniform size
statement. Its proof also gives constant-free ordered arithmetic
formulas over the integers. See Section~\ref{sec:few-order} and
Theorem~\ref{thm:few-order-skolem}.

\paragraph{Result 2: an exact qualitative quantifier-elimination threshold.}
There is a polynomial-time constructible sequence of unambiguous
existential 3-CNFs $\exists z\,\Phi_i(u,z)$ with $N_i$ free bits and
$O(N_i)$ total variables and clauses. Let $S_i(k)$ be the minimum
size after eliminating $z$ into an XOR of at most $k$ layered ordered
parity branching programs. Then, for every sequence $k_i\ge1$,
\[
 \log_2 S_i(k_i)=o(N_i)\quad\Longleftrightarrow\quad k_i\to\infty.
\]
The quantitative bounds are $2^{\lfloor cN_i/2^k\rfloor}$ from below
and $2^{O(N_i/\sqrt{\min\{k,\log N_i\}})}$ from above.
They do not match quantitatively. The family already has linear-size
Boolean circuits; the lower bound concerns its representation after
quantifier elimination. See Corollaries~\ref{cor:diverging-orders-parity}
and~\ref{cor:unambiguous-qe}.

\paragraph{Result 3: restricted structure for all 3-CNF IPS certificates.}
For any diverging order budget, every unsatisfiable $n$-variable 3-CNF
has an integral, constant-free, syntactically multilinear Hilbert-like
IPS certificate of total size $2^{o(n)}$ using at most that many
orders. Its formal degree is at most $n$, and every gate has degree
at most one in the clause placeholders. The order list can be fixed
in advance. The uniform explicit clause-sensitive bound remains
$2^{O((n+m)/\sqrt{\log(n+m+2)})}$. See
Theorems~\ref{thm:few-order-ips} and~\ref{thm:universal-variable-order-ips}.

\paragraph{Result 4: program equivalence and a named proof calculus.}
For ordinary binary regular expressions of full encoding length $n$,
universality, nonuniversality, equivalence, inequivalence, inclusion,
and noninclusion are outside $\mathrm{NTIME}_{\rm mt}[t]$ for every
$t(n)=o(n\sqrt{\log n})$. For every fixed $0<\epsilon<1/2$, fully
written tree-like Kleene-algebra proofs of all valid equations cannot
have worst-case total symbol size
$O(n(\log n)^{1/2-\epsilon})$. This is a superlinear worst-case
bound; we do not exhibit an explicit family of exponentially hard
identities. See Theorems~\ref{thm:regex-nondet-lower}
and~\ref{thm:ka-tree-superlinear}.

\paragraph{Further results.}
Section~\ref{sec:skolem} shows that, for a circuit with at most one
witness per input, the acceptance bit and every witness bit have
few-order representations whose total size obeys the bound of Result~1;
when every input has a witness, these are Skolem functions.
Section~\ref{sec:bounded-variable} separates a bounded-variable
expansion-counting logic from vectorial partial fixed points of the
same arity. Section~\ref{sec:few-order-frege-milestone} states a
reflection hypothesis for the compiled certificates, which we do not
prove, and shows that it would give $2^{o(n+m)}$-size Frege
refutations of 3-CNFs; it also explains why known IPS-to-Frege
simulations do not supply this hypothesis.

\subsection{Dependencies and status}

The results above are consequences of the manuscript~\cite{Cruz26} and
of the explicitly cited published inputs listed below. The simulation
and contraction identities of~\cite{Cruz26} are used as input theorems
and are not reproved here. The new structural arguments, reductions,
size transfers, and the lower-bound specialization are proved in full.
The Frege-reflection hypothesis of Section~\ref{sec:few-order-frege-milestone}
is explicitly conditional. No conjectural class separation is claimed
as a theorem. Our literature searches did not establish priority for
every consequence; the one relevant source we could not consult in full
is identified in Section~\ref{sec:ka-proof-lower}.

We refer to results of~\cite{Cruz26} by descriptive names: its
\emph{simulation theorem} and \emph{counting-classes theorem} (labelled
\texttt{thm:main} and \texttt{thm:counting-classes} in its source), its
\emph{initial-clique lemma} (\texttt{lem:top}, proved in its appendix),
its \emph{merge lemma} (\texttt{lem:merge}), its \emph{formula
construction} (\texttt{sec:formula}), and its \emph{path-index lemma}
(\texttt{lem:paths}). The inputs are used in the following roles.
\begin{enumerate}
\item The simulation theorem of~\cite{Cruz26} proves, for every function
$t(n)\ge n$ and without a constructibility assumption,
\[
 \mathrm{NTIME}_{\rm mt}[t]\subseteq
 \mathrm{DSPACE}[O(t/\sqrt{\log t})].
\]
The counting-classes theorem of~\cite{Cruz26} gives the analogous
inclusions for the probabilistic, parity, fixed-modulus and
equality-count acceptance criteria used below.
\item The initial-clique lemma, the merge lemma, the formula
construction, and the path-index lemma of~\cite{Cruz26} provide the
clique recursion and its potential and path-index bounds. The initial
clique uses $q=2^h$ nodes and has potential
$O((M+q\Delta+q^3)^2)$ for an $M$-edge network of maximum degree
$\Delta$. The proof uses only $q^3\le s$ until the final choice
$h=\Theta(\log s)$.
\item The integral weighted lift of the preceding
note~\cite{PreviousNote} uses the integer coefficient identities
of~\cite{Cruz26} before any numeric specialization. Give every symbolic
weight block one owner vertex. A product of table entries combines
disjoint groups of original vertices, so every product has disjoint
owner scopes. Routing coefficients, tests, grades and split indices
are independent of the symbolic weights. Thus the expanded integer
formula is syntactically multilinear in those blocks. A vertex entry
which is $0,1$, or a linear form in its owner block adds only bounded
local formula size. A root-to-leaf sum of branch-index lengths
$O(F)$ bounds the expanded tree size by $2^{O(F)}$: identify a leaf
by its sequence of child indices, including delimiters charged to
the node overhead in the path-index lemma. Different leaves have
different addresses of $O(F)$ bits, hence there are at most $2^{O(F)}$
leaves and nodes. These observations form the weighted-formula
interface used here; they do not convert a general space algorithm
into a multilinear formula.
\item Disjoint counting sparsification \cite[Lemma~4]{Ant26}, also
stated as assumption (A1) in \cite[Section~2]{Szeider26}: for every
fixed $\epsilon>0$, a 3-CNF on a specified universe of $n$ variables
has its model set partitioned into the model sets of at most
$2^{\epsilon n}$ 3-CNFs, each with at most $M_\epsilon n$ clauses.
This is used only for nonuniform variable-only size bounds. We do not
infer a uniform sublinear-space implementation of sparsification.
\end{enumerate}
Further published inputs are used locally: fast codes with fast duals
\cite{BR26,BIR26} and Homeister's lower-bound method \cite{Homeister03}
in Section~\ref{sec:parity-findings}, the invalid-computation reduction
of Gramlich and Schnitger \cite{GS05} in Section~\ref{sec:regex-lower},
and Kozen's completeness theorem \cite{Kozen94} in
Section~\ref{sec:ka-proof-lower}.

\subsection{Conventions}
\label{sec:conventions}

All logarithms are base two unless a source explicitly uses another
base. For a positive integer $N$, $[N]=\{1,\ldots,N\}$. Machine space counts work-tape cells; inputs are read only.
Every machine has a fixed finite number of tapes and a fixed finite
alphabet. Circuit size includes input vertices, so witness bits are
not free. Arithmetic formula size counts gates and leaves. A
constant-free formula uses only the constants $0,1,-1$. Arithmetic
branching programs may have field constants produced by summing
constant paths; the constant-free assertion applies to the formulas.
Parity programs are over $\mathbb F_2$, and their size counts vertices
and edges.

\paragraph{Ordered formulas.}
The ordered model (Definition~\ref{def:ordered-formula}) constrains
products, not sums: a product must place all variable blocks of its
left factor before all variable blocks of its right factor. Different
branches of a sum may mention the same variables. Over $\mathbb F_2$
addition is XOR, and a parity program accepts exactly when its number
of active accepting paths is odd. These choices permit cancellation
throughout the upper and lower bounds.

\paragraph{Ordered parity programs.}
An ordered parity branching program is layered: each variable
has one layer, all variable layers follow one fixed permutation, and
edges in a variable layer are labelled by $0$, $1$, or $\star$.
An edge with Boolean label $b$ is active when that layer's input bit
is $b$; a $\star$ edge is always active. Additional constant layers
are permitted between variable layers. Multiple edges with the same
label and parity nondeterminism are allowed. Its value is the parity
of the number of active source-to-accepting-terminal paths.
Its size counts vertices and edges, including terminals.
A sum of $k$ such programs is the XOR of their computed functions; the
orders may differ. Its total size is the sum of component sizes.


\subsection{Organization}

Section~\ref{sec:few-order} proves the few-order structure of the
contraction formula, shows how to prescribe the orders in advance, and
makes the number of orders a free parameter.
Section~\ref{sec:ips-universal} derives weighted-model formulas and IPS
certificates with universal order lists, and Section~\ref{sec:skolem}
derives Skolem functions for unique witnesses.
Section~\ref{sec:parity-findings} proves the matching qualitative lower
bound from fast codes and explains why the existing lower bounds do not
transfer to general circuits.
Sections~\ref{sec:regex-lower} and~\ref{sec:ka-proof-lower} prove the
regular-expression and Kleene-algebra lower bounds, and
Section~\ref{sec:bounded-variable} the same-arity separation.
Section~\ref{sec:few-order-frege-milestone} treats the conditional
Frege transfer, and Section~\ref{sec:open-map} lists open problems.

Several small instances are worked out along the way: a one-bit
crossing, a $2\times2$ grid, all four $h=4$ cut histories, a
four-clause unsatisfiable formula, prefix-block examples, a seven-bit
code, a single AND-gate quantifier elimination, a three-cell
transition, and a two-node Kleene-algebra proof. They illustrate the
exact definitions; none replaces a general proof.

\section{Logarithmically many variable orders suffice}
\label{sec:few-order}

This section proves the structural consequence of the integral
contraction construction on which the rest of the paper rests. It uses
the weighted-network lift and formula-size bound of the preceding
note~\cite{PreviousNote}, together with the actual placement of the
sums in the formula construction of~\cite{Cruz26}. The point is that
the clique-value recursion chooses one \emph{global} row or column cut
per merge. All these choices precede the final two-node product, after
which table entries are unfolded, and no independently chosen cuts are
introduced during table unfolding.

Let $s$ be the network size parameter, and let $h$ be the initial name
length, a power of two with $q=2^h$ initial clique nodes. In the choice
made in~\cite{Cruz26}, $h$ is the largest power of two satisfying
$2^{3h}\le s$, and the arithmetic formula has size
$S=2^{O(s/\sqrt{\log s})}$. Original symbolic weights occur at unique
owner vertices. All routing coefficients and consistency tests are
constant with respect to these weights.

\subsection{Cut histories and ordered components}

We first define the class of formulas that the compiler produces.

\begin{definition}[Ordered formulas]
\label{def:ordered-formula}
Fix a linear order $\pi$ of variable blocks. A formula is $\pi$-ordered
if its leaves contain at most one variable block, and at each product
gate every block used by the left child precedes every block used by the
right child. A child with no variable blocks imposes no restriction.
Addition gates may have arbitrary overlaps. Products of more than two
factors are bracketed without changing their displayed order.
\end{definition}

\begin{lemma}[The order determined by a cut history]
\label{lem:cut-history-order}
Put $d=\log_2 h$. Each global cut history
$\sigma\in\{R,C\}^d$ determines a linear order $\pi_\sigma$ of the
initial clique nodes, independent of every grade, split, label, and
table-entry choice. For that history, all products in the unfolded
formula combine consecutive node groups in this order.
\end{lemma}

\begin{proof}
Identify the positions of an original node name with the ordered list
$I=(1,\ldots,h)$. Define a permutation $P_\sigma(I)$ recursively.
If $|I|=1$, let $P_\varnothing(I)=I$. Otherwise split $I$ into its
first and second halves $I_L,I_R$. Set
\[
 P_{R\tau}(I)=P_\tau(I_L)\,I_R,
 \qquad
 P_{C\tau}(I)=P_\tau(I_R)\,I_L,
\]
where juxtaposition concatenates lists. Order original clique nodes
lexicographically by their name bits read in the order $P_\sigma(I)$.
This is $\pi_\sigma$.

At the first row merge, a new node is indexed by the bits $I_L$ and
its table is a sum of products over the old nodes indexed by $I_R$,
taken in lexicographic order. Thus the eventual order of the retained
bits is followed by the unchanged lexicographic order of the discarded
bits. This gives precisely the first recursion above. At a column merge,
the same argument exchanges the two halves and gives the second
recursion. Apply this reasoning at every subsequent merge. At name
length one, the product of the two final tables takes bit values zero
and one in that order, which is the base case.

Consequently, the value expansion orders all original table factors
by $\pi_\sigma$. More specifically, a current table represents the
original nodes having a fixed setting of the retained coordinates.
These coordinates form an initial segment in $P_\sigma(I)$, so its
group is an interval of the final lexicographic node order. Its product
of predecessor tables lists the consecutive subintervals in order.
The intervening factors are routing coefficients or consistency tests
and contain no symbolic weights. Sums over labels and splits change
which entries are selected but not the order of these groups.
\end{proof}

\paragraph{Small order check.}
For $h=2$, the two orders of a $2\times2$ grid are row-major
$00,01,10,11$ and column-major $00,10,01,11$. For $h=4$, the four
bit-position orders are $(1,2,3,4)$, $(2,1,3,4)$,
$(3,4,1,2)$, and $(4,3,1,2)$. These result respectively from
$RR,RC,CR,CC$. No grade index changes an order after its two cut choices
have been fixed.

The next lemma splits the formula into one component per cut history.

\begin{lemma}[Filtering histories without enlarging a component]
\label{lem:filter-histories}
The weighted contraction polynomial computed by the formula
of~\cite{Cruz26} can be written
\[
 W=\sum_{\sigma\in\{R,C\}^{\log_2h}} A_\sigma,
\]
where each $A_\sigma$ is an ordered, syntactically multilinear formula
of size at most $S$, after an inessential constant-factor binary
expansion. The order may depend on $\sigma$, but there are at most $h$
orders in total.
\end{lemma}

\begin{proof}
The clique-value part of the formula consists of nested sums over full
grade-and-split lists, one list per merge. Once its grades are fixed,
the orientation of that entire clique is a determined function of its
budgets and those grades. For a fixed $\sigma$, replace by zero every
child of one of these sum gates whose chosen orientation disagrees
with the corresponding symbol of $\sigma$. Retain all other gates.
Replacing a subtree by a constant leaf cannot enlarge the tree.

Every complete grade-and-split history has one orientation string, so
its contribution survives in exactly one component. Distribution of
the clique-value sums therefore gives the displayed equality.
Contributions from invalid encodings are harmless: their value is
zero in each component where they remain.

Choose once and for all an order of the original vertices within each
initial clique group. In its initial table entry, put the vertex
factors in that order. The coefficient ring is commutative, so this
does not alter the polynomial. Read the weight owners in this order
inside each group, and concatenate these lists in the node order
$\pi_\sigma$ of Lemma~\ref{lem:cut-history-order}. This gives one
global order of the weight blocks. Lemma~\ref{lem:cut-history-order}
shows that all later table products respect this order. The constant-only coefficient
recursions introduce no weight blocks. Thus $A_\sigma$ is ordered.
The unique-owner scope invariant gives syntactic multilinearity.
There are $2^{\log_2h}=h$ histories.
\end{proof}

\begin{lemma}[Ordered formulas to oblivious branching programs]
\label{lem:ordered-formula-roabp}
Let $A$ be a $\pi$-ordered arithmetic formula of size $S$ in $N$
variable blocks, whose nonconstant leaves are linear forms in one
block. Suppose a source--target path in its formula gadget uses at
most one nonconstant edge from each block. It can be converted to an
arithmetic oblivious branching program respecting $\pi$, of size
$O((N+1)S^2)$, with entries affine in one block per layer. In particular,
scalar blocks give a read-once oblivious arithmetic branching program.
For a paired block $\{w_{i,0},w_{i,1}\}$, the matrix of that block
may instead be refined into two scalar layers with a factor-two width
increase.
\end{lemma}

\begin{proof}
For each subformula construct a directed acyclic source--target
gadget. A leaf is one edge; an affine label is split into parallel
constant and homogeneous linear edges. At an addition gate, take disjoint child
gadgets in parallel, with constant-one entrance and exit edges. At a
product, join the first target to the second source in series. Path
weights sum at an addition and multiply at a product, so the resulting
gadget computes $A$ and has $O(S)$ edges. The ordered condition implies
that nonconstant labels on any path occur in $\pi$ order.

To make the conversion explicit, number the gadget vertices
in topological order and let $v=O(S)$ be their number. Let $M$ be
its matrix of constant-edge weights and let $E_j$ be its matrix of
edges carrying a nonconstant linear form from the $j$th block.
All matrices are $v\times v$.
Since $M$ is strictly upper triangular, $M^v=0$. Set
\[
 D=I+M+\cdots+M^{v-1}.
\]
Its entry $D_{u,u'}$ is the sum of weights of constant-only paths from
$u$ to $u'$, including the empty path when $u=u'$. If $a,b$ are the
source and target, the polynomial of the gadget is
\[
 e_a^{\mathsf T}\left(\prod_{j=1}^{N}(I+D E_j)\right)D e_b.
\]
To verify the equality, expand the product. A term selects an increasing
sequence of block positions. Each selected factor chooses a
constant-only path followed by one edge carrying that block; the
last $D$ chooses a constant-only suffix. Every original source--target
path has exactly this decomposition into its nonconstant edges and
maximal intervening constant segments. Conversely, every concatenation
counted by this expansion is a path of the original graph. Orderedness
ensures its nonconstant edges have distinct increasing positions.

The displayed product is an ROABP of width at most $v$, with at most
$O(Nv^2)$ edges, and constant source and target vectors. Every entry
of $I+D E_j$ is affine in the $j$th block. Constant endpoint vectors
can be incorporated into the first or last matrix. The case $N=0$
is a scalar constant. This proves the claimed size bound.

For a weight pair, write its matrix as
$A_0+A_1w_{i,0}+A_2w_{i,1}$. The identity
\[
 A_0+A_1w_{i,0}+A_2w_{i,1}
 =\begin{bmatrix}A_0+A_1w_{i,0}&A_2\end{bmatrix}
  \begin{bmatrix}I\\w_{i,1}I\end{bmatrix}
\]
refines it into two consecutive scalar layers. The intermediate width
doubles and the entrance and exit widths are unchanged.
\end{proof}

\subsection{The compiler and prescribed orders}

Combining the three lemmas gives the compiler for weighted model
counts.

\begin{theorem}[Few-order weighted-count compiler]
\label{thm:few-order-compiler}
Every 3-CNF with $n$ specified variables and $m$ clauses has a weighted
model polynomial that is a sum of $O(\log(n+m+2))$ ordered multilinear
arithmetic formulas, each with constants $0,1,-1$, of total size
\[
 2^{O((n+m)/\sqrt{\log(n+m+2)})}.
\]
Consequently it is a sum of that many ROABPs, with the same asymptotic
total size over any field. More precisely, the number of orders is at
most the initial name length $h$ of the construction in~\cite{Cruz26}.
\end{theorem}

\begin{proof}
The unique-owner weighted network has size $O(n+m+2)$. Apply
Lemma~\ref{lem:filter-histories}; multiplying the formula-size bound
by $h=O(\log(n+m+2))$ preserves its displayed rate. Apply
Lemma~\ref{lem:ordered-formula-roabp} for the final assertion. The
formula identities are integral and hence specialize to any field.
\end{proof}

The orders in Theorem~\ref{thm:few-order-compiler} depend on how the
initial clique groups the owner vertices. The next lemma fixes that
grouping in advance, so that the order list depends only on the number
of owners and on $h$.

\begin{lemma}[Prescribing the owner groups]
\label{lem:prescribed-owners}
Let a network have $M$ edges, maximum degree $\Delta=O(1)$, and
$N\ge1$ symbolic owner vertices of degrees between one and two.
Fix their ordering $1,\ldots,N$. For any initial name length $h$
allowed by the contraction construction, the initial groups can be
chosen so that every owner order $\pi_\sigma$ depends only on $N,h$
and $\sigma$, independently of the rest of the network. The initial
potential remains
\[
 O\bigl((M+2^h\Delta+2^{3h})^2\bigr).
\]
The grouping remains logspace constructible in the full encoding length
of the explicitly given network.
\end{lemma}

\begin{proof}
Put $q=2^h$ and $K=q/2$. As in the proof of the initial-clique lemma
of~\cite{Cruz26}, the first name bit distinguishes
$K$ vertex groups from $K$ test groups. Put owner $i$ in vertex group
\[
 g(i)=\left\lfloor\frac{(i-1)K}{N}\right\rfloor.
\]
Each group contains at most $\lceil N/K\rceil$ owners. To verify
this, its owner indices satisfy
$jN/K\le i-1<(j+1)N/K$, an interval containing at most
$\lceil N/K\rceil$ integers.

Now list the nonowner vertices separately. Let $P$ be their total
number of ports, and put $D=\max(1,\lceil P/K\rceil)$. Assign a
nonowner $v$ to
$\min(K-1,\lfloor\operatorname{pre}(v)/D\rfloor)$, where
$\operatorname{pre}(v)$ counts only ports of preceding nonowners.
The same prefix-count argument as in the initial-clique lemma bounds each
group's nonowner ports by $D+\Delta$: the first such vertex starts
at or after the group's left threshold, and the last crosses its
right threshold by at most its own degree. The last group has at
most $D$ ports. Thus a group's total port load $d_j$ satisfies
\[
 d_j\le 2\lceil N/K\rceil+D+\Delta=O(M/K+\Delta+1),
\]
because $N\le2M$ and $P\le2M$.

Subdivide every network edge and distribute its equality-test vertex
cyclically among test groups, separately for each unordered pair of
vertex groups, exactly as in the initial-clique lemma. Its width estimate uses
only the port load, not the original prefix assignment. Consequently
each vertex--test pair has width at most
$d_j/K+K+1=O(M/K^2+(\Delta+1)/K+K)$.
There are at most $K^2$ nonzero budgets, and the potential estimate
in that proof becomes
$O((M+q\Delta+q^3)^2)$.

Inside each vertex group list owners in increasing owner index and
place all nonowners anywhere after them. Nonowners and test groups
carry no symbolic variables. Therefore restricting any node order
of Lemma~\ref{lem:cut-history-order} to the owners gives a variable
order determined solely by $N,h,\sigma$. Owner groups use division
of $O(\log(M+N+2))$-bit integers. The nonowner prefix counts and the
cyclic test assignments use a constant number of nested scans and
counters, as in the initial-clique lemma. Vertex names use logarithmic space
in the full encoding length; for the intended CNF networks the number
of vertices is polynomial in $M+N$. Hence the construction retains its
logspace bound in the stated input model.
\end{proof}

\paragraph{Prescribed-group check.}
For $N=5$ owners and $K=2$ vertex groups, the groups are
$\{1,2,3\}$ and $\{4,5\}$, with loads at most
$2\lceil5/2\rceil=6$ before nonowners are inserted. The owner
subsequence within either group is unchanged if the network is
replaced by a different sparsification branch.

\begin{corollary}[Few-order Boolean compilation]
\label{cor:few-order-boolean}
Every bounded-fan-in Boolean circuit with $n$ inputs and $g$ gates,
with $s=n+g+2$, is an XOR of $O(\log s)$ ordered read-once parity
branching programs of total size $2^{O(s/\sqrt{\log s})}$.
The orders may differ between summands.
\end{corollary}

\begin{proof}
Take the parsimonious Tseitin 3-CNF of the circuit with an asserted
output. Set both weights of every auxiliary gate variable to one
already at the network stage. Only the original input pairs need
symbolic owner vertices; all other vertices are nonowners. By unique
extension the network polynomial is the model polynomial of the
original Boolean function. Apply the weighted construction and, when
universal orders are required, prescribe these input-owner groups by
Lemma~\ref{lem:prescribed-owners}. Thus the order list depends only on
the number of original inputs and $h$, not the circuit. Reduce
modulo two and set $w_{i,1}=x_i$, $w_{i,0}=1+x_i$ in the
\emph{block} matrix construction of
Lemma~\ref{lem:ordered-formula-roabp}. Each weight pair occupies one
matrix layer, whose entries become affine in the single variable
$x_i$. Thus this produces a read-once program in the induced order
of the original inputs. Its labels are
$0,1,x_i,\neg x_i$ after constant-factor subdivision if needed. The
sum becomes XOR. Substitution and deleting auxiliary positions do not
increase size beyond a polynomial factor.
\end{proof}

\subsection{The number of orders as a parameter}

The maximal choice of $h$ in~\cite{Cruz26} is needed only for the final
size bound. Any smaller admissible $h$ trades size for fewer orders.

\begin{lemma}[Order budget as a parameter]
\label{lem:variable-order-budget}
The few-order compiler is valid for every power-of-two name length
$h\ge1$ with $2^{3h}\le s$, with at most $h$ orders and total size
$2^{O(s/\sqrt h)}$. Its variable-order list may be fixed in advance
from the number of owners and $h$.
\end{lemma}

\begin{proof}
The initial-potential proof needs $2^{3h}\le s$ to give potential
$O(s^2)$, and does not need maximality of $h$. The path-index lemma
of~\cite{Cruz26} bounds the main index contribution by the square
root of that potential divided by $\sqrt h$, namely $O(s/\sqrt h)$.
The explicit routing and budget-bookkeeping terms are at most
$O(2^{2h}\log s\log(h+2)+\log^2s)$. They are
$O(s^{2/3}\log s\log\log(s+2)+\log^2s)
 =o(s/\sqrt{\log s})$, hence are absorbed uniformly for
$1\le h\le(\log_2s)/3$. There remain exactly
$2^{\log_2h}=h$ histories. Multiplying by $h$ and making the
polynomial overhead for branching programs preserves the stated
rate. Apply Lemma~\ref{lem:prescribed-owners} for the last assertion.
\end{proof}

\begin{proposition}[Parameterized compiler]
\label{prop:parameterized-parity-compiler}
Lemma~\ref{lem:variable-order-budget} has the following Boolean form:
for each power of two $h\ge1$ satisfying
$2^{3h}\le s$, a circuit of size $s$ has a sum of at most $h$ ordered
parity branching programs with total size $2^{O(s/\sqrt h)}$.
The implicit constant is independent of $h$ in this range.
\end{proposition}
\begin{proof}
The initial-clique lemma of~\cite{Cruz26}
uses $q=2^h$ groups and bounds its potential by
$O((M+q\Delta+q^3)^2)$.
Here $M=O(s)$, $\Delta=O(\log s)$, and $q^3\le s$, so
$q\Delta=O(s^{1/3}\log s)=O(s)$ and the potential is $O(s^2)$
uniformly in the displayed range of $h$. (The bounded-degree networks
of Lemma~\ref{lem:prescribed-owners} have $\Delta=O(1)$; only the
consequence $q\Delta=O(s)$ is used here.)
Maximality of $h$ is unnecessary for this bound.
The potential-halving proof then gives the scale $F=s/\sqrt h$.
The grade-and-split overhead is
$O(2^{2h}\log s\log(h+2))$, as in Lemma~\ref{lem:variable-order-budget};
dividing it by $F$ gives at most
$O(s^{-1/3}(\log s)^{3/2}\log\log(s+2))=o(1)$.
The table-index estimate and coefficient-recursion estimate in the
path-index lemma of~\cite{Cruz26} are both $O(F)$ with constants
independent of $h$, and the $O(\log^2s)$ charge is also $o(F)$.
The expanded formula therefore has size $2^{O(F)}$.
There are $\log_2h$ global orientation choices, hence at most $h$
orientation strings, and Lemma~\ref{lem:filter-histories} puts each
string into one ordered component.
\end{proof}

\paragraph{Cancellation between order classes is necessary.}
The few-order decomposition cannot be verified by checking each
component separately for zero. In the one-bit crossing of~\cite{Cruz26},
at $(x,y,u,v)=(0,1,0,1)$ the three grades have values
$\Pi_2=-1$, $\Pi_{-2}=-1$, and $\Pi_0=2$. A budget-minimizing merge
uses columns for grade $2$ and rows for grade $-2$. Over
$\mathbb F_2$, the grade-zero contribution vanishes and the two
different-order contributions both equal one. Their sum is zero.
This is already a cancellation between order classes on one crossing.

\paragraph{Limitations.}
The few-order structure does not by itself prove a general Boolean
circuit lower bound. Such a lower bound needs a lower bound for the
\emph{sum} of the allowed ordered programs; a lower bound for a single
order does not suffice. Existing results must be
compared in terms of the total number of Boolean variables, their field,
the permitted number of summands, and the width allowed per summand.
Nor does the representation directly give short Frege proofs. The
known white-box PIT for sums of constantly many differently ordered
ROABPs has double-exponential dependence on the number of summands
\cite[Theorem~3.2]{GKST15}, and this dependence cannot be treated as
a constant when the number of summands grows like $\log s$.


\section{IPS certificates and weighted counts with universal orders}
\label{sec:ips-universal}

We now apply the compiler to refutations. The first theorem compiles a
single network whose value is the first-failed-clause certificate
(its IPS identity follows Grochow and Pitassi~\cite{GP18}). Its size
is governed by the number of variables and clauses. The second and third theorems use disjoint
counting sparsification and the prescribed owner groups of
Lemma~\ref{lem:prescribed-owners} to obtain bounds in the number of
variables alone, with one order list shared by all formulas on a
given universe.

\subsection{A few-order certificate from one network}
\begin{theorem}[Few-order Hilbert-like IPS]
\label{thm:few-order-ips}
Let $F=\bigwedge_{j=1}^m\kappa_j$ be an unsatisfiable 3-CNF with
$n\ge1$ variables and nonempty, nontautological clauses. Write
$\eta_i(0)=1-x_i$, $\eta_i(1)=x_i$, and let
$F_j(\mathbf x)$ be the product of the factors corresponding to the
unique falsifying values of the variables of $\kappa_j$.
There is a constant-free syntactically multilinear formula
$C(\mathbf x,\mathbf y)$ of size
\[
 2^{O((n+m)/\sqrt{\log(n+m+2)})}
\]
such that $C(\mathbf x,\mathbf0)=0$ and
$C(\mathbf x,F_1,\ldots,F_m)=1$ over $\mathbb Z$.
Moreover, $C$ is a sum of $O(\log(n+m+2))$ ordered multilinear
formulas, or a sum of that many ROABPs with the same asymptotic size.
The formula can be chosen to have formal degree at most $n$ and
placeholder degree at most one at every gate.
\end{theorem}

\begin{proof}
Build a network whose Boolean indices include an assignment
$\alpha\in\{0,1\}^n$, clause-truth bits $b_j$, prefix flags $p_j$,
and first-failure flags $z_j$. Impose
\[
 p_0=1,\qquad b_j=\kappa_j(\alpha),\qquad
 p_j=p_{j-1}\wedge b_j,\qquad
 z_j=p_{j-1}\wedge\neg b_j,
\]
and impose $p_m=0$. For each variable $i$, compute
\[
 r_i=\bigvee_{j:i\in\operatorname{vars}(\kappa_j)}z_j,
\]
where an empty disjunction is zero. These constraints have a unique
extension for each assignment $\alpha$. They satisfy the last assertion
precisely when $\alpha$ falsifies at least one clause. In that case
exactly one $z_j$ equals one, and that $j$ is its first falsified clause.

Add one symbolic owner table for each $x_i$, indexed by
$(\alpha_i,r_i)$, whose value is $1$ when $r_i=1$ and
$\eta_i(\alpha_i)$ when $r_i=0$. Add one symbolic owner table for
each $y_j$, indexed by $z_j$, whose value is $y_j$ for $z_j=1$ and
$1$ for $z_j=0$.

The number of Boolean-gate constraints is $O(n+m)$. Each clause has
width at most three, so its truth function uses a constant number of
fan-in-two gates. Let $d_i$ be the number of clauses $\kappa_j$ with
$i\in\operatorname{vars}(\kappa_j)$, that is, the number of inputs
$z_j$ of the disjunction defining $r_i$. The disjunction trees defining
the $r_i$ use $O(\sum_i d_i+n)=O(n+m)$ gates, because $\sum_i d_i$ is
at most the total number of literal occurrences, which is at most $3m$. Each gate relation is ternary. Replace
fan-out by degree-three equality trees; their total number of edges
and vertices is linear in the total number of wire incidences.
The two-index $x_i$ tables and one-index $y_j$ tables already have
degree at most two. We obtain an $O(n+m)$-edge network of maximum
degree three. Every indeterminate has exactly one owner table.

For an assignment with first failed clause $j$, the product of the
symbolic table entries is
\[
 y_j\prod_{i\notin\operatorname{vars}(\kappa_j)}\eta_i(\alpha_i).
\]
All consistency-test entries equal one on its unique extension.
Thus the network value is exactly the first-falsified-clause
certificate
\[
 C(\mathbf x,\mathbf y)
 =\sum_jy_j\sum_{\alpha:\,j\text{ is first false}}
      \prod_{i\notin\operatorname{vars}(\kappa_j)}\eta_i(\alpha_i).
\]
Its two IPS identities follow as in~\cite{PreviousNote}: after
substituting $y_j=F_j$, the omitted factors are restored, and the
assignment sets partition the Boolean cube because $F$ is
unsatisfiable. Their sum is
$\prod_i((1-x_i)+x_i)=1$.

Apply the integral contraction formula construction directly to this
single network. Its scalar table entries are $0,1,x_i,1-x_i,y_j$;
each has a constant-free formula of bounded size. The owner-scope
induction proves syntactic multilinearity in all $n+m$
indeterminates, and initially gives formal degree at most $n+m$.
The explicit displayed certificate has semantic degree at most $n$
and is linear in the $y_j$'s. First apply
Lemmas~\ref{lem:cut-history-order}--\ref{lem:ordered-formula-roabp}
to this one network. Its total number of cut histories is logarithmic
in $n+m+2$, independently of the number of clauses. This gives the
sum of ordered formulas.

We now remove the possible extraneous intermediate degrees without
losing the size bound or any order. The path-index lemma
of~\cite{Cruz26} bounds the sum of product-index lengths along a
path by $O(\log(n+m+2))$. After replacing indexed products by balanced
binary products, the product depth $H$ is therefore
$O(\log(n+m+2))$. For each component $A_\sigma$, keep only its
bihomogeneous terms of $\mathbf y$-degree one and $\mathbf x$-degree
at most $n-1$. To implement this projection as a formula, associate
with each gate its components indexed by $(a,b)$, with
$0\le a\le n-1$ and $b\in\{0,1\}$. At a sum, add corresponding
components. At a product $UV$, use
\[
 (UV)_{a,b}=\sum_{a_1+a_2=a,\ b_1+b_2=b}
                 U_{a_1,b_1}V_{a_2,b_2}.
\]
All degrees are nonnegative, so discarding higher components cannot
affect these identities. For each fixed output degree, a product
duplicates any child component at most $2n$ times. Along any path
there are at most $H$ products. Including the final sum over $a$,
the whole projected tree therefore has at most $(2n)^{H+1}$ times
as many nodes as the original tree. This factor is
$2^{O(\log^2(n+m+2))}$ and is absorbed by the displayed size bound.
Every retained product still combines disjoint variable scopes in
their previous order. Each retained gate has its advertised pair of
degrees, so the final formula has formal degree at most $n$ and
placeholder degree at most one.

The projection is linear, and the total certificate already has
exactly placeholder degree one and $\mathbf x$-degree at most $n-1$:
each nonempty failed clause omits at least one variable. Consequently
the sum of the projected components is still $C$. Applying
Lemma~\ref{lem:ordered-formula-roabp} after projection gives the
claimed ROABPs.
\end{proof}

\paragraph{Four-clause check.}
For $(x\vee z)\wedge(x\vee\neg z)\wedge
(\neg x\vee z)\wedge(\neg x\vee\neg z)$, the four assignments
have first-failure flags $1000,0100,0010,0001$, respectively.
For each assignment both original variables occur in the failed
clause, so both owner tables output one. The network value is
$y_1+y_2+y_3+y_4$, agreeing with the direct IPS calculation.

\subsection{Universal orders and variable-only bounds}

The bound of Theorem~\ref{thm:few-order-ips} depends on the number $m$
of clauses. Disjoint counting sparsification removes this dependence,
at the cost of uniformity. Because the owner groups are prescribed by
Lemma~\ref{lem:prescribed-owners}, all sparsification branches use
the same order list, so components with the same history can be added.

\begin{theorem}[Universal orders for arbitrary 3-CNFs]
\label{thm:universal-variable-order-wmc}
Fix any integer-valued function $h(n)$ taking power-of-two values,
with $h(n)\to\infty$ and $2^{3h(n)}\le n$ for all sufficiently
large $n$. There is, for each $n$, a list $\Pi_n$ of $h(n)$ orders
of the $n$ weight pairs, depending only on $n$ and $h(n)$, with the
following property. The weighted model polynomial of every 3-CNF
on those $n$ specified variables is a sum of $h(n)$ formulas, one
ordered by each member of $\Pi_n$, with worst-case total size
$2^{o(n)}$. The formulas are constant-free and syntactically
multilinear over $\mathbb Z$. The same assertion holds for sums of
ROABPs over any field.
\end{theorem}

\begin{proof}
Fix $\epsilon>0$. The disjoint counting sparsification lemma
\cite[Lemma~4]{Ant26} (see also \cite[Section~2, assumption~(A1)]{Szeider26})
gives
at most $2^{\epsilon n}$ 3-CNFs $G_a$, each with at most
$M_\epsilon n$ clauses, whose model sets partition the model set of
the given formula. The universe of specified variables is kept equal
to $\{1,\ldots,n\}$, including variables unused in a branch.
Therefore $W_F=\sum_a W_{G_a}$ as integral polynomials.

Compile each branch using the same initial name length $h(n)$.
Its network has $O_\epsilon(n)$ edges, and its owner grouping is
prescribed by $n$ and $h(n)$ using
Lemma~\ref{lem:prescribed-owners}. Thus its initial potential is
$O_\epsilon(n^2)$, and the size bound of Lemma~\ref{lem:variable-order-budget} gives
$2^{O_\epsilon(n/\sqrt{h(n)})}=2^{o(n)}$ for each branch.
Most importantly, the list of orders is the \emph{same} list
$\Pi_n$ for every branch and every $\epsilon$.

For each history $\sigma$, add all the branch formulas having that
history. An addition of formulas ordered by one order is still
ordered by that order. This produces $h(n)$ components, of total
size at most $2^{\epsilon n+o(n)}$. If $B(n)$ denotes the largest
minimum size needed over all 3-CNFs on this fixed universe and this
fixed list, then $\limsup_{n\to\infty}\log_2 B(n)/n\le\epsilon$.
Since this holds for every positive $\epsilon$, the limsup is zero.
The integral, constant-free, and scope properties hold for each
branch and are preserved by addition. The branching-program
conversion has polynomial overhead, preserving $2^{o(n)}$.
This is a nonuniform size-existence result; no effective uniform
choice of a decreasing sparsification parameter is asserted.
\end{proof}

\begin{theorem}[Variable-only few-order IPS]
\label{thm:universal-variable-order-ips}
Use the function $h(n)$ and list $\Pi_n$ of the preceding theorem.
Every unsatisfiable 3-CNF on $n$ variables, after removing duplicate
and tautological clauses, has a Hilbert-like IPS certificate over
$\mathbb Z$ which is a sum of $h(n)$ ordered, syntactically
multilinear formulas of worst-case total size $2^{o(n)}$.
The common lists of orders put all clause placeholders first in
their index order, followed by one member of $\Pi_n$ on the original
variables. The certificate has formal degree at most $n$ and
placeholder degree at most one at every gate. The same size bound
holds for the sum-of-ROABPs representation over any field.
\end{theorem}

\begin{proof}
An empty clause has certificate equal to its placeholder, so assume
all clauses are nonempty. Removing duplicates and tautologies leaves
at most $\sum_{j=1}^3 2^j\binom nj=O(n^3)$ clauses. Let $S_j$
be the variables in clause $j$. Form the 3-CNF $G_j$ consisting of
the preceding clauses and the at most three unit clauses that force
the assignment on $S_j$ falsifying clause $j$. Its models are
precisely the assignments whose first false original clause is $j$.

Apply Theorem~\ref{thm:universal-variable-order-wmc} to each $G_j$
on the full universe of $n$ variables. In its formula put both
weights of a variable in $S_j$ equal to one. For $i\notin S_j$,
substitute $w_{i,0}=1-x_i$ and $w_{i,1}=x_i$. The resulting sum
computes
\[
 Q_j(\mathbf x)=
 \sum_{\alpha:\ j\text{ first false}}
             \prod_{i\notin S_j}\eta_i(\alpha_i).
\]
These substitutions preserve each order as a full order with some
positions unused. They preserve syntactic multilinearity because a
product previously combined disjoint weight pairs. Every resulting
formula has formal degree at most $n-|S_j|\le n-1$.

Write $Q_j=\sum_\sigma Q_{j,\sigma}$ with the shared order list.
Set
\[
 C_\sigma(\mathbf x,\mathbf y)=\sum_{j=1}^m
                 y_j Q_{j,\sigma}(\mathbf x),\qquad
 C=\sum_\sigma C_\sigma.
\]
Every displayed product respects the order putting all placeholders
before the $x$-variables. Addition imposes no scope-order restriction,
so $C_\sigma$ is ordered. No product combines two placeholders, and
its formal degree is at most $n$. Multiplication by $y_j$ combines
disjoint variable scopes. The factor $m=O(n^3)$ in total size is
absorbed by $2^{o(n)}$.

After substituting $y_j=F_j$, each assignment's missing factors on
$S_j$ are restored. Since the formula is unsatisfiable, the
first-false sets partition all assignments. Thus
\[
 C(\mathbf x,F_1,\ldots,F_m)
 =\sum_{\alpha\in\{0,1\}^n}\prod_i\eta_i(\alpha_i)
 =\prod_i((1-x_i)+x_i)=1.
\]
Every term contains one placeholder, giving $C(\mathbf x,0)=0$.
All identities have integer coefficients, so they specialize to
every field. The ROABP conversion is polynomial in formula size and
in $n+m=O(n^3)$ and preserves the variable-only size bound.
\end{proof}

\paragraph{Universal-order check.}
For the four-clause
example above each coefficient $Q_j$ is one, so the universal-order
construction returns $C=\sum_j y_j$ in every allowed order list.

\section{Unique witnesses and explicit Skolem functions}
\label{sec:skolem}

The weighted few-order compiler also eliminates existentially
quantified witnesses when they are unique. The resulting representation
theorem applies even when no small deterministic circuit for the
existentially defined function is given. It is not a claim that
first-order logic over finite structures has the Beth property.

\begin{theorem}[Few-order elimination and witness extraction]
\label{thm:few-order-skolem}
Let $V(u,z)$ be a Boolean circuit over a fixed finite basis of fan-in
at most two, with public inputs
$u\in\{0,1\}^n$ and witness bits $z\in\{0,1\}^w$. Assume that for
each $u$ at most one $z$ satisfies $V(u,z)=1$. Let $s$ count all
public and witness input vertices and all gates, plus two.
For any integer $k\ge1$, the functions
\[
 A(u)=[\exists z\ V(u,z)=1],\qquad
 Z_j(u)=[\exists z\ (V(u,z)=1\ \wedge\ z_j=1)]
 \quad(1\le j\le w)
\]
each have an XOR of at most $k$ layered ordered parity programs,
and the total size of all these representations is
\[
 2^{O\!\left(s/\sqrt{\min\{k,\log(s+2)\}}\right)}.
\]
All outputs use the same order list, fixed from $n$ and the chosen
recursion parameter, independently of $V$. If $A(u)=1$, the vector
$(Z_1(u),\ldots,Z_w(u))$ is its unique witness; when $A(u)=0$ the
vector is zero. If every public input has a witness, these are total
Skolem functions for $V$.
\end{theorem}

\begin{proof}
If $n=0$, all the outputs are constants and have constant-size
representations, so assume $n\ge1$. Any other fixed bounded-fan-in
basis can be converted to the stated one with a constant-factor
size increase by expanding each of its finitely many truth tables.
Tseitin-encode $V$ with its output asserted. For each public input
and witness assignment, the gate clauses force exactly one extension
to all gate variables, and that extension satisfies the asserted
output precisely when $V$ accepts. The number of clauses and
variables is $O(s)$.

Build the corresponding weighted network. Give each public bit
$u_i$ one owner with weights $w_{i,0},w_{i,1}$; give all witness
and gate bits weight one. Its polynomial is
\[
 W(w)=\sum_{a\in\{0,1\}^n}
      \#\{z:V(a,z)=1\}\prod_{i=1}^n w_{i,a_i}.
\]
For the $j$th output use the same circuit constraints and additionally
assert $z_j=1$, yielding
\[
 W_j(w)=\sum_{a\in\{0,1\}^n}
      \#\{z:V(a,z)=1,\ z_j=1\}\prod_{i=1}^n w_{i,a_i}.
\]
Both networks have $O(s)$ edges and bounded degree after equality
trees replace fan-out. Their symbolic owners are the same $n$
public inputs. Prescribe their initial groups by
Lemma~\ref{lem:prescribed-owners}, using a common power of two $h$.
All their components consequently respect the same history-order list.

Substitute $w_{i,0}=1+u_i$ and $w_{i,1}=u_i$ over $\mathbb F_2$.
On a Boolean assignment $u$, the product indexed by $a$ is one
exactly when $a=u$ and zero otherwise. Thus the resulting outputs
are the parity of the displayed witness counts. The promise makes
each count zero or one. The first is therefore $A(u)$, and the
$j$th is $Z_j(u)$.

For sufficiently large $s$, choose the largest power of two at most
$\min\{k,(\log_2s)/3\}$. Then $h\le k$, $2^{3h}\le s$, and
$h=\Theta(\min\{k,\log s\})$. Apply
Lemma~\ref{lem:variable-order-budget} and the block-matrix conversion
in Lemma~\ref{lem:ordered-formula-roabp}; literal pinning uses one
matrix layer per public input. The bound for each output is
$2^{O(s/\sqrt h)}$. There are at most $s+1$ outputs, so their total
size has logarithm $O(s/\sqrt h)+O(\log s)=O(s/\sqrt h)$.
The finitely many smaller $s$ can use full assignment expansions
in one fixed order, after increasing the absolute constant.
Uniqueness proves that the simultaneous bits are the actual witness
when one exists, and their definitions give zero otherwise.
\end{proof}

\paragraph{Hand checks.}
For $V(x,z)=[z=x]$, every $x$ has one witness. The two polynomials
are $W=w_{1,0}+w_{1,1}$ and $W_1=w_{1,1}$. Pinning gives
$A(x)=1$ and $Z_1(x)=x$. In contrast, if $V(x,z)=1$ for both values
of one witness bit, then $W=2(w_{1,0}+w_{1,1})$ reduces to zero
modulo two, although existential acceptance is one. The unique-witness
promise cannot be omitted. For acceptance alone it may be weakened
to the promise that every nonzero witness count is odd; that weaker
promise does not by itself justify extracting a consistent witness
bit by bit. For example, a relation accepting exactly the witnesses
$00,01,10$ has odd total count, but the two bitwise parities give
$11$, which is not an accepted witness.

\section{Fast codes and the exact order threshold}
\label{sec:parity-findings}

This section shows that the few-order upper bound is qualitatively
tight. The code lower bound below does not use the counting simulation.
It is a self-contained specialization of the mechanism of
Homeister~\cite[Theorem~1]{Homeister03}, and we make no novelty claim
for the method. The representation upper bound follows from the
few-order compiler of Section~\ref{sec:few-order}. The combined result
is a consequence of that compiler and of published 2026 code
constructions~\cite{BR26,BIR26}; it is not a general Boolean circuit
lower bound.

\subsection{A lower bound for codes}

We use the layered model of ordered parity branching programs from
Section~\ref{sec:conventions}. All vector spaces in this section are
over $\mathbb F_2$.

The lower bound rests on a combinatorial lemma about prefixes of $k$
orders, a rank bound for a diagonal tensor, and the standard extension
property of codes with large dual distance.

\begin{lemma}[Prefix blocks]
\label{lem:prefix-blocks-parity}
Let $\pi_1,\ldots,\pi_k$ be permutations of a set of $n$ variables,
and let $1\le \ell\le n$.  Put $b=\lfloor\ell/2^k\rfloor$.
If $b\ge1$, there are prefixes $P_i$ of $\pi_i$ and sets $B_i\subseteq P_i$
such that
\begin{enumerate}
\item $|B_i|=b$ for every $i$;
\item any two $B_i$ are equal or disjoint;
\item $B_j\subseteq P_i$ if $B_j=B_i$, and
      $B_j\cap P_i=\varnothing$ otherwise;
\item $|P_1\cup\cdots\cup P_k|<\ell$.
\end{enumerate}
\end{lemma}
\begin{proof}
For $i\ge1$, put $h_i=\lfloor\ell/2^i\rfloor$.
We build prefixes and provisional blocks through stage $i$ so that
properties 2 and 3 hold for the first $i$ indices, each provisional
block has size at least $h_i$, and the union $U_i$ of the first $i$
prefixes has size at most $\sum_{j=1}^{i}h_j$.
For stage 1, take the first $h_1$ variables in $\pi_1$ as both $P_1$
and $B_1$.

Suppose stages through $i-1$ have been completed.
Scan $\pi_i$ until its scanned prefix $P$ satisfies one of the following:
$|P\setminus U_{i-1}|=h_i$, or
$|P\cap D|=h_i$ for a distinct provisional block $D$.
Such a prefix exists because every old block has at least
$h_{i-1}\ge2h_i$ variables.  At the first stopping time just one
condition becomes true: a newly scanned variable lies either outside
$U_{i-1}$ or in at most one of its distinct old blocks.
All other counters are at most $h_i-1$.
Set $P_i=P$.

If $|P\setminus U_{i-1}|=h_i$, set $B_i=P\setminus U_{i-1}$ and
replace each old $B_j$ by $B_j\setminus P$.
Each old block retains at least
$h_{i-1}-(h_i-1)\ge h_i$ variables.
The new block is outside every old prefix; old blocks retain their
containment or disjointness relations with all old prefixes.
Only the new block meets the new prefix.
Thus properties 2 and 3 are preserved.

Otherwise there is one distinct old block $D$ with $|P\cap D|=h_i$.
Set $B_i=P\cap D$.  Replace every old block equal to $D$ by this same
$B_i$, and replace every other old block $B_j$ by $B_j\setminus P$.
The latter blocks retain at least $h_i$ variables.
For any old prefix $P_m$, the block $D$ was either contained in $P_m$
or disjoint from $P_m$.  The first alternative held precisely when
the old block $B_m$ was $D$.  Consequently the new block $B_i$ has
exactly the required relation with $P_m$.
The other blocks preserve all old relations and avoid $P_i$.
Properties 2 and 3 again hold.

In both cases $|P_i\setminus U_{i-1}|\le h_i$, so the bound on $|U_i|$
is preserved.  After stage $k$, shrink each distinct block to any
$b$-element subset, making the same choice for all copies of that
block.  The three block properties are preserved.  Finally,
$|U_k|\le\sum_{i=1}^{k}\lfloor\ell/2^i\rfloor<\ell$.
\end{proof}

\paragraph{Hand checks for Lemma~\ref{lem:prefix-blocks-parity}.}
For $k=2$ and $\ell=8$, let $P_1=\{1,2,3,4\}$ and let the second order
start $5,1,6,2,\ldots$.  The second scan stops at
$P_2=\{5,1,6\}$.  Its new block is $\{5,6\}$; the old block becomes
$\{2,3,4\}$ and may be shrunk to $\{2,3\}$.
The union of prefixes has size 6, and each prefix contains only its
own surviving block.
For the other stopping case, start the second order $3,5,1,6,\ldots$.
It stops at $\{3,5,1\}$; both blocks become $\{1,3\}$, and the union
has size 5.
For $k=3$ and $\ell=16$, take the first prefix $\{1,\ldots,8\}$ and
let the second order start $9,1,10,2,11,3,12,4,\ldots$.
The first two blocks become $\{4,5,6,7,8\}$ and $\{9,10,11,12\}$.
Let the third order start $4,9,13,10,\ldots$.
The third prefix is $\{4,9,13,10\}$; the second and third blocks both
become $\{9,10\}$, while the first block may be shrunk to $\{5,6\}$.
The prefix union is $\{1,\ldots,13\}$, of size $13<16$.

\begin{lemma}[A diagonal tensor with one protected factor]
\label{lem:protected-diagonal-parity}
Let $D\ge1$ and $q\ge1$.  Define a tensor on $[D]^{q+1}$ by
\[
 T(a_1,\ldots,a_q,t)=\prod_{j=1}^{q}[a_j=t].
\]
Suppose
\[
 T=\sum_{j=1}^{q}\sum_{v=1}^{r_j}
       u_{jv}(a_j)\,v_{jv}(a_1,\ldots,a_{j-1},a_{j+1},\ldots,a_q,t).
\]
Then $\sum_jr_j\ge D$.
\end{lemma}
\begin{proof}
Let $U_j\le\mathbb F_2^D$ be the span of the vectors $u_{jv}$.
It has dimension at most $r_j$, so it contains at most $r_j$ of the
coordinate unit vectors $e_1,\ldots,e_D$: any collection of these
vectors is linearly independent.
If $\sum_jr_j<D$, choose $t_0$ such that $e_{t_0}\notin U_j$ for
every $j$.  There is a linear functional vanishing on $U_j$ but taking
value 1 on $e_{t_0}$: extend a basis of $U_j$ by $e_{t_0}$ and then
to a basis of the whole space, and prescribe its values on that basis.
Write that functional as inner product with $\lambda_j\in\mathbb F_2^D$.
Thus $\lambda_j(t_0)=1$ and
$\sum_a\lambda_j(a)u_{jv}(a)=0$ for every $v$.

Contract the displayed tensor identity against
$\prod_j\lambda_j(a_j)$, summing over every $a_j$.
Every term on its right is zero because the contraction in its
designated coordinate kills $u_{jv}$.
Its left becomes the vector $t\mapsto\prod_j\lambda_j(t)$, which
takes value 1 at $t_0$.  This contradiction proves the claim.
\end{proof}

\paragraph{Hand check for Lemma~\ref{lem:protected-diagonal-parity}.}
For $D=2,q=2$, the only nonzero entries of the tensor are
$T(1,1,1)=T(2,2,2)=1$.
If $r_1+r_2<2$, at least one coordinate vector is absent from both
designated spans.  Contracting at that coordinate leaves its diagonal
entry equal to 1 and kills every proposed slice.
For $D=4$, the proposed spans
$U_1=\langle e_1\rangle$ and
$U_2=\langle e_2,e_3\rangle$ are killed by contraction with $e_4$
in both coordinates, whereas the diagonal entry indexed by 4 survives.

\begin{lemma}[Distance implies extensions]
\label{lem:code-extensions-parity}
Let $C\le\mathbb F_2^n$ have dual distance $d^\perp$.
For any $U\subseteq[n]$ with $|U|<d^\perp$, every assignment to $U$
extends to a codeword of $C$.
\end{lemma}
\begin{proof}
If the image of the linear projection $C\to\mathbb F_2^U$ were a
proper subspace, extending a basis would give a nonzero linear
functional on $\mathbb F_2^U$ vanishing on that image.
Extend its coefficient vector by zeros outside $U$.
The resulting vector belongs to $C^\perp$, is nonzero, and has weight
at most $|U|<d^\perp$.  This contradicts the definition of dual distance.
\end{proof}

\begin{theorem}[Code membership requires many orders or large size]
\label{thm:code-orders-parity}
Let $C\le\mathbb F_2^n$ have minimum distance $d$ and dual minimum
distance $d^\perp$.  Put $\ell=\min\{d,d^\perp\}-1$.
If its characteristic function is the XOR of $k$ ordered parity
branching programs of total size $S$, then
\[
             S\ \ge\ 2^{\lfloor\ell/2^k\rfloor}.
\]
\end{theorem}
\begin{proof}
If $\lfloor\ell/2^k\rfloor=0$, the assertion is $S\ge1$ and follows
because the characteristic function accepts the zero codeword.
Otherwise apply Lemma~\ref{lem:prefix-blocks-parity} to the $k$ orders.
Write $U=\bigcup_iP_i$ and enumerate the distinct surviving blocks as
$D_1,\ldots,D_q$, each of size $b=\lfloor\ell/2^k\rfloor$.
Fix every variable in $U\setminus\bigcup_jD_j$ to 0.

For each $t\in\{0,1\}^{b}$, assign $t$ to every $D_j$, using a fixed
enumeration of each block.  By Lemma~\ref{lem:code-extensions-parity},
this assignment to $U$ extends to a codeword; choose one extension
$z(t)$ on $[n]\setminus U$.
If instead the block assignments are $a_1,\ldots,a_q$, while the
outside assignment remains $z(t)$, the resulting word belongs to $C$
exactly when all $a_j=t$.
One direction follows from the choice of $z(t)$.
For the other, a different accepted word would differ from this
codeword only within $U$, and hence have Hamming distance at most
$|U|<\ell<d$, a contradiction.
The evaluations of the characteristic function therefore form the
tensor of Lemma~\ref{lem:protected-diagonal-parity} with $D=2^b$.

For the $i$th program, cut its layering just after the last variable
of $P_i$, placing any subsequent constant layers on the suffix side.
Every accepting path crosses exactly one vertex in the cut layer.
Decomposing paths by that vertex writes the
computed function as
\[
                  \sum_{v} A_v(x_{P_i})B_v(x_{[n]\setminus P_i}),
\]
where the number of summands is at most the program's size.
The formula counts prefix and suffix paths modulo 2; concatenation
at their common cut vertex is a bijection with accepting paths.
After our restriction, $A_v$ depends on only the one residual block
contained in $P_i$.  Substitution of the outside assignment $z(t)$
leaves the other factor depending on the other blocks and $t$.
Thus the $k$ programs give a slice representation of the protected
diagonal tensor with at most $S$ slices.  The preceding lemma yields
$S\ge2^b$.
\end{proof}

\paragraph{Hand check for the code step.}
Take the binary Hamming code of length 7 with parity-check columns
$001,010,011,100,101,110,111$.  It has distance 3: no one column is
zero, no two columns coincide, and the first three sum to zero.
Every nonzero combination of the three check rows has weight 4, so
the dual distance is 4.  Hence $\ell=2$, and with one order the block
size in the proof is 1.  Take $U=\{1\}$.
The assignments $x_1=0$ and $x_1=1$ extend to $0000000$ and
$1110000$.  Using outside strings $000000$ and $110000$, respectively,
the four evaluations form the $2\times2$ identity matrix: the two
off-diagonal words have syndromes $001$.  This checks both extension
and uniqueness in the smallest example used here.

\paragraph{Relation to Homeister.}
Theorem~\ref{thm:code-orders-parity} is a fixed-order, layered
specialization with no graph-ordering denominator: its prefix cuts
require no pigeonhole choice of graph-ordering nodes.
Homeister~\cite[Theorem~1]{Homeister03} proves
$2^{\Omega(\min\{d,d^\perp\}/2^k)}/\prod_i|G_i|$ for graph-driven
components. The proof above is a specialization of that method, and we
claim no new lower-bound technique.

\subsection{An explicit separation from fast codes}

To apply Theorem~\ref{thm:code-orders-parity} to a function with small
circuits, we need codes whose membership test is cheap. Transposition
reduces this to a fast encoder for the dual code.

\begin{lemma}[Transposing a linear encoder]
\label{lem:transpose-parity}
If an $n$-output linear map $u\mapsto Mu$ over $\mathbb F_2$ has a
fan-in-two XOR circuit with $a$ gates, then $x\mapsto M^\top x$ has
an XOR circuit of size $O(a+n)$.  Consequently, if a code's dual
has such an encoder, membership in the code has a Boolean circuit
of size $O(a+n)$.
\end{lemma}
\begin{proof}
Attach a formal output variable $x_j$ to the vertex producing output
$j$.  Process the original circuit backwards.  At a gate $v=p+q$,
send the accumulated value at $v$ to both $p$ and $q$; whenever a
vertex receives several values, combine them by binary XOR gates.
At an original input the accumulated value is the corresponding
coordinate of $M^\top x$.
Indeed the coefficient of $x_j$ at input $u_i$ is the parity of the
number of paths from $u_i$ to the original output $j$, which is the
coefficient $M_{ji}$ obtained by expanding the original XOR circuit.
There are $2a$ gate-input edges and $n$ output attachments, so all
accumulations use $O(a+n)$ gates.
If the columns of $M$ span $C^\perp$, then $x\in C$ exactly when
$M^\top x=0$.  Testing this with a binary OR tree and one negation
adds $O(n+a)$ gates.
\end{proof}

\paragraph{Hand check for transposition.}
For $M(u_1,u_2)=(u_1,u_2,u_1+u_2)$, one XOR gate suffices.
Reversing it yields $M^\top(x_1,x_2,x_3)=(x_1+x_3,x_2+x_3)$.
The common contribution $x_3$ is sent to both original inputs.

\paragraph{Fast codes with fast duals.}
Brehm and Resch~\cite[Theorem~2]{BR26} give binary rate-$1/2$ codes
$C,C^\perp$ with both distances $\Omega(n)$ and both encoders of
linear size.  Their construction is randomized; it does not directly
give a uniform deterministic choice of a good code.
Brehm, Ishai and Resch~\cite[Theorem~9]{BIR26} strengthen this to
polynomial-time-uniform randomized systematic encoders with $O(n)$
random bits, $O(n)$ wires, and both distances $\Omega(n)$ except with
probability $2^{1-\Omega(n)}$.
Section~2.2 of~\cite{BIR26} defines uniformity as polynomial-time
construction of circuit descriptions and defines size as the number of
wires, and Remark~3 there states logarithmic depth.
The formal theorem is expressed on an increasing sequence of block
lengths $n_i$; the conclusions below are stated on that sequence.

\begin{theorem}[An explicit representation separation]
\label{thm:fast-code-order-separation}
Use the few-order compiler of Corollary~\ref{cor:few-order-boolean}.
There is a polynomial-time-uniform family of functions $F_i$ on
$N_i=\Theta(n_i)$ bits, each having a Boolean circuit with $O(N_i)$
gates and depth $O(\log N_i)$, such that
\begin{enumerate}
\item every representation of $F_i$ as the XOR of $k$ ordered parity
      branching programs has size at least
      $2^{\lfloor cN_i/2^k\rfloor}$, for one absolute $c>0$;
\item $F_i$ has such a representation using $O(\log N_i)$ orders
      and total size $2^{O(N_i/\sqrt{\log N_i})}$.
\end{enumerate}
For each fixed $k$, the lower bound is exponential in $N_i$.
For any fixed $K>0$, a representation of size at most
$2^{K N_i/\sqrt{\log N_i}}$ must use at least
$\tfrac12\log_2\log_2N_i-O_K(1)$ orders.
\end{theorem}
\begin{proof}
Use \cite[Theorem~9]{BIR26} with $q=2$, limiting rate $1/2$,
and any fixed $\epsilon\in(0,1/4)$.
Let $E_i(r,m)=(m,P_i(r,m))$ be its systematic encoder, with $r$
denoting its $O(n_i)$ random bits, now treated as ordinary input bits.
For $x=(m,y)\in\{0,1\}^{n_i}$ define
\[
                    F_i(r,m,y)=[y=P_i(r,m)].
\]
Replacing field additions and multiplications by XOR and AND gates,
respectively, gives an $O(n_i)$-size circuit for $P_i$.
Bitwise comparison followed by a balanced AND tree adds $O(n_i)$
gates and $O(\log n_i)$ depth.
The circuit family remains polynomial-time uniform.
Since $n_i\le N_i\le A n_i$ for a constant $A$, its size and depth
have the claimed bounds in $N_i$.

For all sufficiently large $i$, the failure probability for either
distance is less than 1.  Choose a good seed $r_i$ solely for the
lower-bound argument.  Its code and dual have distance at least
$\delta n_i$ for one fixed $\delta>0$.
Restricting the seed bits of any ordered parity branching program
cannot increase its size or create an additional order: delete the
inactive edges in a fixed-bit layer and make its other edges constant.
Keep that layer as a constant layer, retaining path parity.
Thus any $k$-order representation of $F_i$ restricts to one of
the characteristic function of the good code.
Theorem~\ref{thm:code-orders-parity} applies with
$\ell\ge\delta n_i-2\ge cN_i$, after decreasing $c$ and discarding
finitely many initial indices.  This proves the lower bound.

Apply the compiler of Corollary~\ref{cor:few-order-boolean} to the
explicit $O(N_i)$-size circuit to
obtain the upper bound.  If $S\le2^{K N_i/\sqrt{\log N_i}}$, the lower
bound implies
\[
 \frac{cN_i}{2^k}-1\le\log_2S
       \le\frac{K N_i}{\sqrt{\log_2N_i}}.
\]
For all sufficiently large $i$, this yields
$2^k\ge(c/(K+1))\sqrt{\log_2N_i}$, and taking logarithms proves
the final assertion.
\end{proof}

\paragraph{Why the seed trick is legitimate.}
No algorithm for finding a good seed is asserted or needed.
The input function evaluates every seed using one explicit circuit.
The lower bound uses existence of a good restriction, and restrictions
preserve the representation model.  This gives a uniform hard function
from a randomized code construction without derandomizing that
construction.  The uniformity is polynomial-time uniformity, not
DLOGTIME uniformity.

\paragraph{Scope of the separation.}
Theorem~\ref{thm:fast-code-order-separation} shows an
exponential-to-subexponential change when the number of
allowed variable orders grows from any fixed constant to logarithmic.
The current lower and upper order counts for the displayed
subexponential size are $\Omega(\log\log N)$ and $O(\log N)$.
It also refutes a proposed universal lower bound $2^{\Omega(n)}$
against $O(\log n)$ orders based solely on positive rate and linear
primal and dual distance: the fast-code family supplies counterexamples.
It gives no lower bound against unrestricted Boolean circuits;
the exhibited functions already have linear-size circuits.

Combining the code lower bound with the parameterized compiler
(Proposition~\ref{prop:parameterized-parity-compiler}) gives the
qualitative threshold.

\begin{corollary}[Exact qualitative threshold]
\label{cor:diverging-orders-parity}
For the explicit family $F_i$ of
Theorem~\ref{thm:fast-code-order-separation}, let $S_i(k)$ be its
minimum total size as an XOR of at most $k$ ordered parity programs.
For any sequence of positive integers $k_i$,
\[
             \log_2 S_i(k_i)=o(N_i)
             \quad\Longleftrightarrow\quad k_i\longrightarrow\infty.
\]
The forward implication uses only the code lower bound.
The reverse implication holds for every family with linear-size
Boolean circuits under the parameterized compiler.
\end{corollary}
\begin{proof}
If $k_i$ does not tend to infinity, some fixed $K$ satisfies $k_i\le K$
at infinitely many indices.  On those indices the code lower bound
gives $\log_2 S_i(k_i)\ge cN_i/2^K-1$, which is not $o(N_i)$.
If $k_i\to\infty$, let $s_i=O(N_i)$ be the linear-size circuit bound,
including all input vertices.  Choose $h_i$ as the largest power of
two at most $\min\{k_i,(\log_2s_i)/3\}$.
Then $h_i\to\infty$, and the parameterized compiler gives at most
$h_i\le k_i$ orders and log size $O(s_i/\sqrt{h_i})=o(N_i)$.
\end{proof}

\paragraph{Interpretation of the threshold.}
An arbitrarily slowly diverging number of permitted orders suffices
for subexponential size, whereas any bounded number fails on the
explicit fast-code family.  This is an exact qualitative threshold
for this representation model.  It is not an exact quantitative
size--order tradeoff: the lower exponent is proportional to $N/2^k$,
whereas the compiler's upper exponent is $O(N/\sqrt{k})$ until
$k$ reaches order $\log N$.

\begin{corollary}[Unambiguous quantifier elimination]
\label{cor:unambiguous-qe}
There is a polynomial-time constructible family of 3-CNFs
$\Phi_i(u,z)$ with $O(N_i)$ clauses and variables, with $N_i$ free
bits $u$, such that each free assignment has at most one satisfying
extension to $z$, and the relation
\[
                \exists z\,\Phi_i(u,z)
                \ =\ \bigoplus_z\Phi_i(u,z)
\]
has the size--order threshold of
Corollary~\ref{cor:diverging-orders-parity}.
Here the normal form after elimination is an XOR of ordered parity
branching programs; cancellation is permitted both inside and between
its components.
\end{corollary}
\begin{proof}
Take the explicit linear-size circuit for $F_i$ and give each gate
output a fresh witness bit.
For a fan-in-two gate $w=B(a,b)$, forbid each of the four triples
$(a,b,w)$ inconsistent with its Boolean truth table by a width-three
clause.  For a forbidden triple $(\alpha,\beta,\gamma)$, the clause
has literal $a$ when $\alpha=0$ and $\neg a$ when $\alpha=1$, and
the corresponding choices for $b$ and $w$.
This clause is false on exactly that triple.
Constant gates use unit clauses, and one final unit clause requires
the circuit output to be 1.

For each free assignment, processing the circuit in topological order
forces exactly one value on every gate output.  This unique complete
assignment satisfies the final unit clause precisely when the circuit
accepts.  Thus zero or one witness exists, as asserted, and its
existential and parity projections both equal $F_i$.
There are at most four clauses per binary gate, one per constant gate,
and one final clause, so the formula size is $O(N_i)$ and its
construction is polynomial-time uniform.
The preceding representation bounds apply to its projected relation.
\end{proof}

\paragraph{Hand check for the quantifier step.}
For an AND gate $z=x\wedge y$ followed by acceptance of $z$, one may
use the four truth-table exclusion clauses and the unit clause $z$.
The triples $000,010,100,111$ are exactly the four gate-consistent
triples.  The unit clause retains only $111$.
Each free pair $(x,y)$ therefore has one accepting witness exactly
when $x=y=1$, and none otherwise.

\subsection{Why existing lower bounds do not transfer to circuits}

A lower bound for the few-order representation of a function with
small circuits would give a circuit lower bound through the compiler,
so it is important to compare the known lower bounds with the
compiler's parameters. None of the following results is strong
enough.

\paragraph{Set-multilinear branching programs.}
We use the corrected revision~2 of Chatterjee, Kush, Saraf and
Shpilka~\cite{CKSS24}. The original title and abstract incorrectly
identified their set-multilinear model with arbitrary sums of ROABPs;
the revision records the correction.
Theorem~1.6 of~\cite{CKSS24} proves $\exp(\Omega(n^{1/3}))$ total-width lower bounds
for polynomials in $\Theta(n^2)$ variables and degree $\Theta(n)$.
Theorem~1.9 proves $n^{\Omega(d)}$ for a \emph{single} ordered
set-multilinear ABP, not for arbitrary sums.
Lemma~3.1 gives, for $t\le e^{d/96}$ summands of maximum width $w$,
\[
 t w^q\ge n^{\sqrt{dq}/20},
 \qquad q\ge 12\ln(2t\sqrt d),\quad q\le d/4,
\]
with its stated integrality conventions; the alternative in the
lemma is $t>e^{d/96}$.
For $t=O(\log s)$ this gives an exponent roughly
$\sqrt{d/\log d}\log n$, below the $d\log n$ bits needed to index a
coefficient.  It therefore does not provide a superlinear Boolean
circuit lower bound through the compiler.
The arc polynomial used for some statements is analyzed over
characteristic zero; reducing its coefficients modulo 2 cannot be
assumed to preserve the lower bound.

\paragraph{Sums of graph-driven parity programs.}
Homeister's paper~\cite{Homeister03} (journal version in TCS, 2006)
covers sums of graph-driven parity programs, so it
does apply to sums of parity OBDDs. Its dependence on $k$ is
$2^k$ in the denominator of the exponent.
Even with distance $\Omega(n)$, this becomes only
$2^{\Omega(n^{1-c})}$ when $k=c\log_2n$.
This is below the compiler's $2^{O(n/\sqrt{\log n})}$, so the existing
lower bound cannot imply a superlinear general circuit bound.

\paragraph{Identity testing for sums of ROABPs.}
Gurjar, Korwar, Saxena and Thierauf~\cite[Theorem~3.2]{GKST15} give
white-box PIT time $(nd\,w^{2^c})^{O(c)}$ for $c$ summands of width $w$,
$n$ variables, and individual degree $d$.
The dependence on $c$ is not polynomial or merely single exponential.
Putting $c=\Theta(\log s)$ and $\log w=O(s/\sqrt{\log s})$ into this
bound does not yield a polynomial-time verifier in the original size
$s$, or a near-linear-exponent verifier for the compiled certificate.
The compiler therefore does not by itself supply efficient proof
checking for this arithmetic representation.

\paragraph{General multilinear ABPs.}
Bor\'em Fabris, Limaye, Srinivasan and
Yehudayoff~\cite[Question~1 and Theorems~3 and~7]{FLSY26} state the
open problem
of explicit superpolynomial multilinear ABP lower bounds and analyze
limitations of full min-partition rank.  This remains a different
model from a logarithmic sum of ordered ABPs.
Neither a formula lower bound nor a lower bound for one order can
be substituted for the needed few-order lower bound.

\subsection{A conditional coefficient-predicate bridge}

The following proposition makes precise which few-order lower bound
would give a circuit lower bound. It requires an extension of the
compiler to owner vertices of logarithmic arity, which we do not
prove.

\begin{proposition}[Conditional coefficient bridge]
\label{prop:coefficient-bridge}
Suppose the few-order weighted-network compiler also permits a
vertex of arity $O(\log s)$ whose table entries are distinct formal
weight variables owned by that vertex.
Let $g:[N]^d\to\{0,1\}$ have a Boolean circuit of size $s$ on
$d\lceil\log_2N\rceil$ input bits, where $N\le s^{O(1)}$ and inputs
count toward $s$.  Then
\[
 P_g(X)=\sum_{a\in[N]^d}g(a)\prod_{j=1}^{d}X_{j,a_j}
\]
is a sum of $O(\log s)$ ordered set-multilinear ABPs, of total size
$2^{O(s/\sqrt{\log s})}$, over $\mathbb F_2$.
\end{proposition}
\begin{proof}
Give each Boolean circuit wire a Boolean network edge and each gate
its truth-table constraint.  For every input tuple accepted by the
circuit, all internal wire values are uniquely determined.
Enforce acceptance at the output.  For block $j$, place one owner
vertex on its $\lceil\log_2N\rceil$ input bits, with table value
$X_{j,a}$ on a valid encoding of $a\in[N]$ and 0 on unused encodings.
The contraction of this network is exactly $P_g$, because accepted
inputs have one extension and rejected inputs have none.
The network has $O(s)$ edges and vertices.  Each owner has arity
$O(\log s)$ and a polynomial-size table.
The assumed compiler gives the claimed bound; since every product
combines disjoint vertex sets, every monomial uses at most one
variable from each owner block.  Grouping by global orientation
strings preserves this property and gives the asserted ordered
set-multilinear components.
\end{proof}

\paragraph{Hand check for the coefficient bridge.}
For $d=2,N=2$ and $g(a,b)=[a=b]$, the two accepted input tuples
produce exactly
$X_{1,1}X_{2,1}+X_{1,2}X_{2,2}$.
An equality gate has one valid internal extension for each accepted
tuple, so the coefficients are 1, not path counts larger than 1.

\paragraph{The required lower bound.}
If one could prove $N^{\Omega(d)}$ total size for
$O(\log(d\log N))$ different-order components for such a polynomial,
then the coefficient predicate, on $b=d\log N$ Boolean bits,
would require $\Omega(b\sqrt{\log b})$ general Boolean circuit size.
Indeed the lower and upper bounds imply
$s/\sqrt{\log s}=\Omega(b)$.
If $s\ge b^2$ the conclusion already holds; otherwise
$b\le s<b^2$ makes $\log s=\Theta(\log b)$ and yields the claim.
The lower bound of~\cite{CKSS24} falls far short of
$N^{\Omega(d)}$ for a sum of logarithmically many orders.

\section{A nondeterministic lower bound for regular-expression equivalence}
\label{sec:regex-lower}

This section derives a consequence of the nondeterministic and counting
time--space simulations of~\cite{Cruz26}. The invalid-computation
reduction is classical; the new deduction is its superlinear
nondeterministic-time consequence at the scale of the full bit length.
We spell out a linear-size version for ordinary expressions, because
citing PSPACE-completeness or an $O(n)$-state NFA alone would leave a
logarithmic encoding gap. The expressions below use only letters,
$\epsilon$, $\varnothing$, binary union, binary concatenation, and
Kleene star. There is no intersection, complement operator, squaring
operator, binary exponent, or shared subexpression. Every exponent in
the proof abbreviates an explicit list of that many copies.

\paragraph{The published reduction.}
Gramlich and Schnitger~\cite[Lemma~7 and the appendix proof of
Theorem~4]{GS05} construct invalid-computation automata with $O(|x|)$
states from fixed deterministic in-place machines, and explain that
the regular-expression version gives the analogous linear gap. We use
the extended version of their conference paper. The proofs below supply
the details for ordinary expressions and their encoding that we need.

\subsection{A linear-size reduction}

\begin{lemma}[A local configuration model]
\label{lem:regex-local-machine}
For every $L\in\mathrm{DSPACE}[n]$ there is a fixed deterministic
one-tape decider $M$ whose configurations on an input $x$ of length $n$
have a fixed length $S=n+2$, over a fixed finite alphabet $B$, with
exactly one state-marked tape cell. Its initial configuration
$c_0(x)=a_1\cdots a_S$ is obtainable by a fixed letterwise encoding of
$x$ with endpoint markers. There is a fixed finite relation
\[
 T\subseteq (B\cup\{\#\})\times B\times
                  (B\cup\{\#\})\times B
\]
such that, for well-formed configurations $c,c'$ of length $S$, $c'$ is
a one-step successor of $c$ if and only if
\[
 (\hat c_{j-1},\hat c_j,\hat c_{j+1},c'_j)\in T\quad(1\leq j\leq S),
\]
where $\hat c=\#c\#$ is $c$ padded by separators, so that
$\hat c_j=c_j$ for $1\le j\le S$ and $\hat c_0=\hat c_{S+1}=\#$.
No halting configuration has a successor. The separator $\#$ is not
in $B$.
\end{lemma}
\begin{proof}
Fix a finite multitape decider for $L$ using at most $cn$ cells, increasing
$c$ to an integer. Retain the input in a read-only track and represent
all work symbols and virtual head markers in a fixed number of tracks
of length $n$. A cell of the new tape stores the tuple of symbols and
head-marker bits at one position. This uses a fixed finite alphabet.
The one physical head can sweep these tracks to locate the finitely many
virtual heads, retain their scanned symbols and original-machine state
in finite control, and perform a simulated transition in further sweeps.
Crossing from one length-$n$ work segment to the next is implemented
by a sweep to the opposite endpoint and a finite-control change of
segment. The simulation's time is immaterial here. Put fixed endpoint
markers in two additional cells. It never moves outside these $S=n+2$
cells and halts exactly when the original decider halts. Its initial
tracks are the letterwise input encoding, blanks, and designated head
markers.

Encode a physical-head cell holding tape symbol $b$ and state $q$ by
$(q,b)$, and other cells by their plain tape symbols. The successor
symbol at a position depends on its three-cell neighborhood. At the
head position it is the newly written symbol, with a state marker
retained only for a stationary move. At a neighboring position it
acquires the new state marker exactly when the head moves into it;
otherwise its symbol is unchanged. A position not adjacent to the head
keeps its symbol. These are finitely many rules, hence give the finite
relation $T$. For a halting state, permit no next symbol for a triple
whose center has that state marker. Reject illegal boundary moves.
On a well-formed nonhalting configuration these rules determine exactly
its machine successor. On a halting configuration its head-position
rule forbids every proposed successor. This proves both directions.
\end{proof}

For a fixed finite set of letters $A$, write $[A]$ for their union as an
ordinary regular expression. This notation never denotes an operator
that complements a regular language. Put $\Gamma=B\cup\{\#\}$.
The set $H\subseteq B$ consists of state-marked letters, and
$H_{\mathrm{acc}}\subseteq H$ consists of those with accepting state.

\begin{lemma}[Linear-length invalid-computation expressions]
\label{lem:linear-regex-reduction}
For each fixed machine $M$ in Lemma~\ref{lem:regex-local-machine}, there
is a deterministic multitape transducer mapping $x$ to an ordinary
regular expression $R_x$ over $\Gamma$ such that
\[
 L(R_x)=\Gamma^*\setminus
 \{\#c_0(x)\#c_1\#\cdots\#c_t\#:
          c_0,\ldots,c_t\text{ is an accepting computation of }M\}.
\]
The output has syntax length $\Theta_M(n+1)$ and is produced in
$O_M(n+1)$ time. Its symbol at a given output position can be recovered
in $O_M(\log(n+2))$ workspace. In particular,
\[
 x\in L(M)\quad\Longleftrightarrow\quad L(R_x)\ne\Gamma^*.
\]
Since $M$ is deterministic and halts, the omitted set is either empty
or a singleton.
\end{lemma}
\begin{proof}
For $n\geq1$ we have $S\geq3$; finitely many smaller inputs can be
handled separately. Form $R_x$ as the union of the following constant
number of expressions. Constants here may depend on $M$.

\emph{Wrong delimiters or block lengths.} Include
\[
 \epsilon\; +\; \#\; +\; [B][\Gamma]^*\; +\;
 [\Gamma]^*[B]\; +\;
 [\Gamma]^*\#(\epsilon+[B])^{S-1}\#[\Gamma]^*
 \; +\;[\Gamma]^*[B]^{S+1}[\Gamma]^*.
\]
The first four alternatives cover the empty word, the single separator,
and words without a leading or trailing separator. For a remaining
word, the blocks between consecutive separators consist of letters of
$B$. The fifth alternative detects a block of length less than $S$,
since $(\epsilon+[B])^{S-1}$ denotes all $B$-words of length at most
$S-1$. The sixth detects a block of length greater than $S$. Therefore
these alternatives accept exactly the words failing the shape
$\#B^S(\#B^S)^*\#$.

\emph{An ill-formed configuration.} Include
\[
 [\Gamma]^*\#[B\setminus H]^*\#[\Gamma]^*
 \quad+\quad
 [\Gamma]^*[H][B]^*[H][\Gamma]^*.
\]
On words of the required shape, the first finds a block without a head
marker and the second finds two head markers in the same block. They
accept no word all of whose configurations have exactly one head.

\emph{A wrong initial configuration.} Write $c_0(x)=a_1\cdots a_S$.
Define an expression recursively by
\[
 D(\epsilon)=\varnothing,\qquad
 D(a_1u)=[B\setminus\{a_1\}][\Gamma]^*+a_1D(u),
\]
and include $\#D(c_0(x))$.
On a word of the required shape, this accepts exactly when the first
configuration differs from $c_0(x)$: choose the first differing
position, match its correct prefix, use the corresponding mismatch
alternative, and accept the remaining suffix. Conversely every match
of this expression identifies a differing position. The recursive
expression has length $O_M(S)$ because $D(u)$ occurs once, not twice.
This factoring is the point that avoids the quadratic union over all
possible mismatch positions.

\emph{A nonaccepting last configuration.} Include
\[
 [\Gamma]^*\#[B\setminus H_{\mathrm{acc}}]^*\#.
\]
Because the expression ends at the end of the word and its last starred
alphabet excludes $\#$, it checks precisely the last block. On a
well-formed block, it accepts exactly when its state is not accepting.

\emph{A bad transition.} For every fixed forbidden quadruple
$(a,b,c,d)\notin T$, with $b,d\in B$ and $a,c\in\Gamma$, include
\[
 [\Gamma]^*abc[\Gamma]^{S-1}d[\Gamma]^*.
\]
In a word of the required shape, a match with center $b$ at a cell of
one configuration places $d$ at the corresponding cell of the next
configuration. Indeed, the start of the next block is $S+1$ positions
later; after $c$, exactly $S-1$ intervening symbols precede $d$.
The letters $a,c$ can be separators at the two boundary cells.
If there is no next block, the required $d\in B$ cannot be matched.
Thus these alternatives detect precisely a position violating the
local transition relation. Lemma~\ref{lem:regex-local-machine} proves
that absence of every such violation is equivalent to a correct step.

Every invalid word violates one of the five requirements. Every valid
accepting computation satisfies all five and is rejected by all
alternatives. This proves the language identity.

There are only $O_M(1)$ forbidden quadruples. Every displayed power is
expanded into that many explicit copies, and every such power has
exponent $O(S)$. The recursively factored initial test has $O(S)$ nodes.
Consequently the full syntax length is $O_M(S)$, and the long-block test
also makes it $\Omega_M(S)$.

A fixed number of left-to-right scans of $x$ produces all repeated
blocks and the initial-configuration test. To serialize the latter
in prefix notation, emit its constant template and current letter,
then the template and next letter, finishing with $\varnothing$.
Fully parenthesized notation is obtained by a further scan emitting
the requisite closing delimiters. All syntax tokens have constant
bit length over the fixed alphabet. Endpoint padding costs a constant.
Thus the transducer takes $O_M(n+1)$ time.
For access to one output symbol, regenerate the output while counting
positions; all loop counters and input positions use $O_M(\log(n+2))$
bits. This gives the required small-space input interface.
\end{proof}

\begin{lemma}[One fixed binary alphabet]
\label{lem:regex-binary}
The reductions in Lemma~\ref{lem:linear-regex-reduction} can all target
ordinary expressions over the single fixed alphabet $\{0,1\}$, with
the same asymptotic length, construction-time, and access-space bounds.
\end{lemma}
\begin{proof}
For the fixed alphabet $\Gamma$ of a particular reduction choose an
injective fixed-length code $h:\Gamma\to\{0,1\}^d$; $d$ is constant
for that reduction. Replacing each letter of $R_x$ by its code gives
an ordinary expression $h(R_x)$ whose language is the encoded language.
Let $A=0+1$, and let $K$ be the finite set of unused $d$-bit codewords.
The following constant-length expression accepts exactly the binary
strings which are not concatenations of valid codewords:
\[
 B_h=(A^d)^*(A+A^2+\cdots+A^{d-1})
       +(A^d)^*[K]A^*.
\]
When $d=1$, the first summand is $\varnothing$. The first summand
checks a nonzero length residue modulo $d$. For lengths divisible by
$d$, the second finds an invalid aligned block; it also harmlessly
accepts some strings with a nonzero residue.
Put $\widehat R_x=B_h+h(R_x)$. If a word is not properly coded,
$B_h$ accepts it. If it is $h(w)$, injectivity of the block code gives
$h(w)\in L(\widehat R_x)$ exactly when $w\in L(R_x)$.
Therefore universality and nonuniversality are preserved, and the
complement is a singleton when the original complement is a singleton.
All added expressions and letter replacements have size bounded by a
constant depending on $M$, so all three bounds are preserved.
\end{proof}

\subsection{The lower bound and its variants}

\begin{theorem}[Natural superlinear nondeterministic lower bound]
\label{thm:regex-nondet-lower}
Assume the simulation theorem of~\cite{Cruz26},
$\mathrm{NTIME}_{\mathrm{mt}}[t]\subseteq
\mathrm{DSPACE}[O(t/\sqrt{\log t})]$ for $t(n)\geq n$.
Let $\mathrm{RE\mbox{-}UNIV}$ be universality of ordinary binary
regular expressions, encoded as syntax trees in prefix or fully
parenthesized notation, and let $\mathrm{RE\mbox{-}NONUNIV}$ be its
complement on valid encodings. Neither problem belongs to
\[
 \mathrm{NTIME}_{\mathrm{mt}}[t(n)]
 \qquad\text{for any }t(n)=o(n\sqrt{\log n}).
\]
The same assertion holds for regular-expression equivalence and inclusion,
and for their complementary decision problems. The size parameter $n$
is full encoding length, not number of states or variable names.
\end{theorem}
\begin{proof}
Every $L\in\mathrm{DSPACE}[n]$ reduces by the preceding lemmas to
$\mathrm{RE\mbox{-}NONUNIV}$ with output length $\Theta_L(n)$ and
deterministic construction time $O_L(n)$. Applying the construction to
a decider for the complement of $L$ gives the same reduction to
$\mathrm{RE\mbox{-}UNIV}$.

Suppose either target had an $o(n\sqrt{\log n})$-time nondeterministic
multitape algorithm. Compose it with the appropriate reduction, storing
the linear-length output on a tape and then running that algorithm.
No monotonicity of its time bound $t$ is needed.  For suitable fixed
constants $C_L,D_L$, a bound on the composed time is
\[
 T_L(n)=D_L(n+1)+D_L\max_{1\le m\le C_L(n+1)}\lceil t(m)\rceil.
\]
This is at least $n$ after increasing $D_L$.
It is also $o(n\sqrt{\log n})$: given $\eta>0$, beyond a fixed
$m_0$ we have $t(m)\le\eta m\sqrt{\log(m+2)}$.
The finitely many smaller values have a fixed maximum, while all
remaining terms in the displayed maximum are bounded by
$\eta C_L(n+1)\sqrt{\log(C_L(n+1)+2)}$.
Dividing by $n\sqrt{\log n}$ and then letting $\eta$ tend to zero
proves the assertion.  This envelope need not be constructible:
the simulation theorem of~\cite{Cruz26} explicitly permits arbitrary
time bounds. That theorem gives a deterministic space bound
\[
 O\!\left(\frac{T_L(n)}{\sqrt{\log T_L(n)}}\right)
 \leq O\!\left(\frac{T_L(n)}{\sqrt{\log n}}\right)=o(n).
\]
This would put every language in $\mathrm{DSPACE}[n]$ in
$\mathrm{DSPACE}[o(n)]$, contradicting the deterministic space hierarchy.
For completeness we give a single-language diagonalization against the
whole union $\mathrm{DSPACE}[o(n)]$. Enumerate one-binary-work-tape deterministic machines with
an endmarker-bounded read-only input tape by finite descriptions
$M_1,M_2,\ldots$, using length-lexicographic enumeration of descriptions
and default machines for malformed descriptions. Description $i$ has
length $O(\log(i+2))$. On unary length $n$, put $i=1+\nu_2(n)$ and
simulate $M_i$ using at most $n$ binary work cells. If it tries to
cross that cap, give the bounded simulation a fixed rejecting outcome.
If $q_i$ is its state count, the number of capped configurations is
at most $q_i(n+2)(n+1)2^{n+O(1)}$. Stop after this many steps and
give a nonhalting simulation the same default outcome. Because
$i\le1+\log_2n$, its description, tape contents, head positions,
and a configuration counter all fit in $O(n)$ space. Define a tally
language to return the opposite of this bounded outcome, and reject
nonunary inputs.

Every fixed multitape $o(n)$-space decider has an equivalent normalized
machine in this enumeration with $o(n)$ space: encode its fixed
alphabet and tapes by constant-length blocks and tracks and simulate
a transition by sweeps. If $M_j$ decided the diagonal language, it
would fit the $n$-cell cap for all sufficiently large lengths with
$1+\nu_2(n)=j$. There are infinitely many such lengths. On each of
them it halts and cannot repeat a configuration before halting, so
the bounded simulation returns its actual answer. The diagonal
language returns the opposite answer, a contradiction. Thus the
language lies in $\mathrm{DSPACE}[n]$ outside the whole union
$\mathrm{DSPACE}[o(n)]$.

Finally, universality of $R$ is its equivalence to $(0+1)^*$, and is
also the inclusion $(0+1)^*\subseteq L(R)$. Pairing $R$ with this
fixed expression increases input length and construction time by a
constant. Nonuniversality gives the complementary decisions.
\end{proof}

The same argument applies to the counting acceptance criteria, using
the counting-classes theorem of~\cite{Cruz26}.

\begin{corollary}[Counting and probabilistic acceptance]
\label{cor:regex-counting}
Under the corresponding counting-class simulations of~\cite{Cruz26},
all the problems in Theorem~\ref{thm:regex-nondet-lower} are outside
$\mathrm{PPTIME}[t]$, $\mathrm{BPTIME}[t]$, $\oplus\mathrm{TIME}[t]$,
$\mathrm{Mod}_q\mathrm{TIME}[t]$ for fixed $q$, and
$\mathrm C_=\mathrm{TIME}[t]$, for every
$t(n)=o(n\sqrt{\log n})$.
\end{corollary}
\begin{proof}
The simulations of~\cite{Cruz26} count the two transition choices as distinct
computations even when their transition functions agree, so the
deterministic prefix must be implemented with its acceptance
interpretation in mind.  For probabilistic machines, identical
transitions in the prefix leave all eventual acceptance probabilities
unchanged, and thus preserve both bounded and unbounded error.
For $\mathrm C_=\mathrm{TIME}$, if the prefix takes $r(x)$ steps,
identical choices multiply the eventual accepting and rejecting
computation counts by the same positive factor $2^{r(x)}$.
Their equality is preserved.

For parity and $\mathrm{Mod}_q\mathrm{TIME}$, instead give each
prefix step one continuing transition and one transition that halts
and rejects immediately.  There is exactly one continuing computation
through the full prefix.  All other prefix computations reject, so
the number of accepting computations is exactly the target machine's
number, with no multiplicative factor.  This preserves its parity
and its residue modulo $q$.
Once the reduction and input-head initialization are complete,
simulate each transition choice of the target machine without adding
further branching steps.

In every case the composition has time
$T_L(n)=o(n\sqrt{\log n})$. Apply the small-space simulation
of~\cite{Cruz26} for the relevant acceptance interpretation and repeat
the space-hierarchy contradiction.
\end{proof}

\begin{corollary}[A verification-work obstruction]
\label{cor:regex-verifier}
For every fixed $0<\epsilon<1/2$, there is no sound and complete
certificate verifier for ordinary binary regular-expression universality
which (i) runs in deterministic multitape time $O(n+|\pi|)$ on an
expression of length $n$ and certificate $\pi$, and (ii) supplies every
universal expression with some certificate of length
$O(n(\log n)^{1/2-\epsilon})$.
\end{corollary}
\begin{proof}
Guess a certificate up to the stated, explicitly clockable bound and
run the verifier. Completeness supplies an accepting branch exactly
on universal expressions, and soundness forbids one otherwise. The
running time is $O(n(\log n)^{1/2-\epsilon})$, which is
$o(n\sqrt{\log n})$, contradicting
Theorem~\ref{thm:regex-nondet-lower}.
\end{proof}

\paragraph{Scope.}
Theorem~\ref{thm:regex-nondet-lower} is an unconditional superlinear
nondeterministic-time lower bound for a familiar formal-language
decision problem, derived from the simulation theorem of~\cite{Cruz26}.
It is not an exponential lower bound and does not separate NP from
PSPACE. Corollary~\ref{cor:regex-verifier} is parameterized by
its stated verification cost; a particular Kleene-algebra proof calculus
must have that cost verified before the corollary is applied to it.
No deterministic linear-space upper bound for NFA universality is used.
The usual nondeterministic subset-search upper bound is linear in the
number of NFA states; invoking Savitch only gives a quadratic
space bound at that parameter.

\paragraph{Hand check of the reduction.}
For $S=3$, take two blocks $abc$ and $def$, so the word is
$\#abc\#def\#$. To inspect the center $b$ and corresponding next
letter $e$, the pattern is $abc[\Gamma]^2e$: the two intervening
characters are $\#$ and $d$. For the first cell, $\#ab[\Gamma]^2d$
has intervening characters $c,\#$; for the last cell,
$bc\#[\Gamma]^2f$ has intervening characters $d,e$.
The short-block expression $(\epsilon+[B])^2$ accepts lengths
zero, one and two, while $[B]^4$ detects lengths at least four.
The initial-prefix test on $abc$ expands as
\[
 [B\setminus\{a\}][\Gamma]^*
 +a\bigl([B\setminus\{b\}][\Gamma]^*
          +b([B\setminus\{c\}][\Gamma]^*+c\varnothing)\bigr),
\]
which rejects a correct $abc$ prefix and accepts each possible first
mismatch in exactly the corresponding branch.

\section{Proof size in tree-like Kleene algebra}
\label{sec:ka-proof-lower}

We now derive a proof-size consequence from the regular-expression
lower bound of Section~\ref{sec:regex-lower} and the simulation
of~\cite{Cruz26}. It is a superlinear symbol-size lower bound for a
precisely specified tree-like presentation of Kleene algebra. It is not
a superpolynomial proof lower bound, and it does not assert the same
verification rate for arbitrary proofs with shared references.

\paragraph{Kozen's axioms.}
We use Kozen's axiomatization~\cite{Kozen94}. Section~2 of the author
version, equations (3)--(17), gives the finite equational and Horn
axioms, and Theorem~19 (completeness) states that equal regular
languages have a Kleene-algebra proof. We checked these statements
directly in the author version.


\begin{definition}[The ground tree presentation of Kleene algebra]
\label{def:tree-ka}
Fix two alphabet letters $a,b$. A term is an ordinary fully written
regular expression over $a,b,0,1$ using the binary operations $+$
and concatenation and unary star. The constants $0,1$ denote the
empty language and the language containing the empty word.
Terms have no sharing, exponent notation, or extension names.
An inequality $u\le v$ abbreviates the equation $u+v=v$.

The initial equations are all ground instances of Kozen's
equations (3)--(15). The two induction rules are the ground instances
\[
 \frac{b+ax\le x}{a^*b\le x},
 \qquad
 \frac{b+xa\le x}{ba^*\le x}.
\]
The letters $a,b,x$ in these two displayed schemas denote term
metavariables, not the fixed alphabet letters. The equality rules
are reflexivity, symmetry, transitivity, and congruence for $+$,
concatenation, and star. Congruence for a binary operation takes two
equality premises, and congruence for star takes one. Every rule is
thus one of finitely many fixed patterns with at most two premises
and a constant number of term metavariables.

A proof is a finite rooted tree whose nodes are full equations,
with a rule name and, if desired, its explicit instantiating terms.
Its size is its total number of symbols over a fixed finite encoding
alphabet. We count all copies of every term and judgment.
\end{definition}

\begin{lemma}[Soundness and completeness of the ground presentation]
\label{lem:tree-ka-complete}
An equation over the fixed alphabet has such a finite proof if and
only if its two regular languages are equal.
\end{lemma}

\begin{proof}
Each semiring axiom holds for regular languages under union and
concatenation. The star-unfolding axioms hold because every word in
$1+AA^*$ or $1+A^*A$ belongs to $A^*$. If $B+AX\subseteq X$,
then $B\subseteq X$ and $AX\subseteq X$. Induction on $k$ gives
$A^kB\subseteq X$: the case $k=0$ is $B\subseteq X$, and the
next case follows by multiplying the preceding inclusion on the
left by $A$. Taking the union over all $k$ gives
$A^*B\subseteq X$. The other induction rule is proved by right
multiplication in the same way. Equality rules preserve equality
of languages. Induction up a finite proof tree proves soundness.

For completeness, let $u\sim v$ mean that the ground equation
$u=v$ has a finite tree proof. Reflexivity, symmetry, transitivity,
and congruence make this a congruence on the term algebra. The
quotient by $\sim$ satisfies all the Kleene-algebra axioms. For
the Horn axioms, if a premise holds in the quotient, its ground
equation has a proof by the definition of $\sim$; placing the
corresponding induction rule above that proof gives a proof of its
conclusion. Thus the quotient is a Kleene algebra. Kozen's
Theorem 19 says that equality of the two regular languages implies
their equality in every Kleene algebra, including this quotient.
Hence $u\sim v$, which by definition gives a finite tree proof.
\end{proof}

\begin{lemma}[Linear-time verification]
\label{lem:tree-ka-linear-check}
There is a deterministic machine with a fixed finite number of
ordinary work tapes that verifies an explicitly annotated tree
proof of size $L$, and checks that its conclusion is a supplied
equation of length $n$, in time $O(L+n)$.
Moreover, any tree proof with full judgments and rule names can be
annotated for this checker with at most a constant-factor increase
in total symbol size.
\end{lemma}

\begin{proof}
Encode each term in prefix notation; each symbol's arity is zero,
one, or two. A term's boundary and syntactic validity can be checked
in linear time with a unary pending-argument counter: start with
one pending position, remove one when reading a symbol and add its
arity, and require the first return to zero to be the end of the
term. This uses a fixed number of tape moves per symbol.

List proof nodes in postorder. Each node contains its rule tag,
the instantiating terms for that fixed rule, and its full conclusion.
The verifier maintains a stack of the conclusions of the already
verified subtrees. On reading a rule with $k\le2$ premises, it pops
the last $k$ conclusions, rejecting if fewer are present. Copy those
conclusions onto $k$ scratch tapes. A stack block can be copied in
reverse while popping and then read from the opposite end; this
costs at most twice its length and needs no random-access operation.
Copy each instantiating term onto
a separate scratch tape; the number of these tapes is bounded by
the maximum number of metavariables in the finite rule list.

For each fixed premise pattern, scan the pattern and the corresponding
popped conclusion simultaneously. At a constructor symbol compare
that symbol. At a metavariable, scan its stored term against the
corresponding segment of the conclusion and then return the term
tape to its left endpoint. Each metavariable occurs only a bounded
number of times in the finite pattern, so the work is linear in the
total length of the local instantiating terms and premise judgments.
Check the node's conclusion against the instantiated conclusion
pattern by the same procedure. Push its verified full conclusion
onto the stack. Clearing and resetting a scratch tape costs at most
its previous contents plus its new contents, so summed over all
nodes these resets preserve the linear total bound.

Each judgment is pushed once and popped at most once because the
proof is a tree. Every full term is scanned a bounded number of
times per local rule use. Stack movement is charged to symbols
written or removed; it never scans earlier unrelated stack blocks.
At the end require exactly one stack entry, compare it with the
supplied equation, and reject any trailing malformed input.
These operations have total time $O(L+n)$ on a fixed finite number
of tapes. Postorder tags also check that the input describes one
tree: each rule removes exactly its arity many completed subtrees
and replaces them by one.

For the size comparison, remove unused metavariables from every
rule schema. Each remaining instantiating term appears at least
once in a premise or conclusion of that local rule. Therefore the
sum of the lengths of a node's instantiating terms is at most the
sum of lengths of its local judgments. Across a tree, each judgment
is charged at most twice, once as its node's conclusion and once
as its parent's premise. The complete annotation has $O(L)$ size.
\end{proof}

\begin{theorem}[Superlinear Kleene-algebra tree proof size]
\label{thm:ka-tree-superlinear}
Assume the simulation theorem of~\cite{Cruz26} and the linear-size,
linear-time reduction of Section~\ref{sec:regex-lower} from
linear-space computations to ordinary binary regex equivalence.
For every fixed $0<\epsilon<1/2$, there
is no constant $c$ such that every valid equation of ordinary binary
regular expressions of total length $n$ has a ground tree
Kleene-algebra proof of at most
\[
 c\,n(\log(n+2))^{1/2-\epsilon}
\]
symbols. More generally, no time-constructible function
$B(n)\ge n$ with $B(n)=o(n\sqrt{\log n})$ can uniformly bound
all these shortest proof sizes up to a constant factor.
\end{theorem}

\begin{proof}
Suppose first that the displayed bound holds. Replace $\epsilon$
by a smaller positive rational number if necessary; the hypothesized
bound then still holds, and its exponent is computable. On input an equation
of length $n$, guess an annotated proof of at most a fixed constant
multiple of that length bound. The constant accounts for the
annotation in Lemma~\ref{lem:tree-ka-linear-check}. Reject a guess
exceeding the bound. Run that verifier and accept precisely if it
accepts. All branches halt within
$O(n(\log(n+2))^{1/2-\epsilon})$ steps. The bound is
time-constructible, so a clock enforces it with the same order of
time. Soundness rejects every invalid equation, and completeness
together with the hypothesized size bound accepts every valid one.
This would put regex equivalence in
\[
 \operatorname{NTIME}
   [O(n(\log(n+2))^{1/2-\epsilon})],
\]
contradicting the lower bound of Theorem~\ref{thm:regex-nondet-lower},
$\mathrm{REGEX\mbox{-}EQUIV}\notin
\operatorname{NTIME}[o(n\sqrt{\log n})]$.
For a time-constructible $B$, use its clock in the identical argument.
\end{proof}

\paragraph{Small-instance check.}
The one-node proof of $(a+b)+0=a+b$ instantiates $x+0=x$
with $x=a+b$. Its annotation and judgment both have constant size.
For a two-node proof, take this judgment as the premise of symmetry
to conclude $a+b=(a+b)+0$. The postorder verifier pushes the first
judgment, pops it at symmetry, compares it with the premise pattern,
and pushes the reversed equation. It never searches a list of earlier
proof lines.

\paragraph{Limits and related work.}
The proof does not establish a superpolynomial lower bound for
Kleene algebra and does not settle $\mathsf{NP}\ne\mathsf{PSPACE}$.
An arbitrary nonconstructible little-$o$ upper envelope of shortest
proof lengths does not immediately furnish a clocked nondeterministic
algorithm; the theorem therefore states the explicit logarithmic
powers and the time-constructible generalization. Reusing previously
proved equations by arbitrary binary line references also requires
a different verification analysis.

The chapter of Beier and Holzer~\cite{BH24} on the proof complexity
of Salomaa's axiom system $F_1$ is directly relevant to novelty. We
checked its publisher table of contents but could not obtain the full
chapter, and no theorem from it is used as an input to a proof. A
reproduced abstract states a tower upper bound for $F_1$ and the
equivalence between polynomially bounded general regex-equivalence
proof systems and $\mathsf{NP}=\mathsf{PSPACE}$; those statements do
not by themselves provide an existing stronger lower bound for the
tree model above. We do not claim a completed literature search
establishing novelty for Theorem~\ref{thm:ka-tree-superlinear}.


\section{A same-arity separation for bounded-variable counting}
\label{sec:bounded-variable}

The following separation is a corollary of the clause-parameter
counting simulation of~\cite{Cruz26}. It applies to arbitrary
relational $k$-variable first-order matrices, not only to noncrossing
certificates. It is not a solution of Fagin's arity hierarchy
problem, which allows arbitrarily many first-order variables. We make
no priority claim for this restricted-fragment separation. In
particular, the $k=2$ existential fragment on successor words was
already shown to be regular by Gr\"adel and Rosen~\cite{GR99}.

\paragraph{Models.}
Fix an integer $k\geq2$. Structures have a universe $[N]$, $N\geq2$,
a fixed finite relational vocabulary of maximum arity $k$, and a linear
order. A word uses unary letter predicates with the order and endpoints
implicit. A relation table can equivalently be supplied explicitly;
its entries can be addressed by tuples in $O(\log N)$ workspace.
No function symbols are allowed. The formula is fixed, not input.
For a sentence $\varphi(\bar R)$ of $\mathrm{FO}^k$, using at most
$k$ variable names but permitting their reuse and arbitrary quantifier
nesting, let $c_\varphi(A)$ be the number of interpretations of a fixed
finite tuple $\bar R$ of additional relations of arity at most $k$
which make $\varphi$ true. Unused relation entries are counted as choices.
Let $\mathcal K_k$ be the class defined by finite Boolean combinations
of comparisons $c_\varphi\leq c_\psi$ and fixed-modulus tests on these
counts. Count operators are not nested in first-order matrices.

The first step grounds a bounded-variable sentence as a 3-CNF whose
size is governed by the arity, not by the length of the formula.

\begin{lemma}[Parsimonious bounded-variable grounding]
\label{lem:fo-grounding}
For every fixed $\varphi$ there is a logspace-accessible $3$-CNF
$G_{A,\varphi}$ with $O_\varphi(N^k)$ variables and clauses such that
\[
 \#G_{A,\varphi}=c_\varphi(A).
\]
Here logspace accessibility means that a clause and one of its literals
can be recovered from their indices in $O_\varphi(\log N)$ workspace
with read-only access to $A$.
\end{lemma}
\begin{proof}
Introduce one Boolean input for every tuple in every quantified relation.
Their number is $O(N^k)$. For each syntactic subformula $\theta$ and
each assignment to its free variables, introduce a gate for its truth
value under that assignment. There are at most $N^k$ gates per
subformula before replacing unbounded fan-in.

A relation atom refers to its corresponding Boolean input or to a
fixed input-structure bit. Equality and order atoms become constants.
At a Boolean connective, the gate refers to its child gates with the
assignment restricted to the child's free variables. At a quantifier
$\exists x\,\theta$ or $\forall x\,\theta$, combine the $N$ child
gates by OR or AND. If $x$ does not occur freely in $\theta$, remove
this vacuous quantifier. Otherwise, if the quantified formula has $r$
free variables, its child has $r+1\leq k$. Binary bracketing of all its
quantifier gates costs $O(N^{r+1})=O(N^k)$. The resulting bounded-fan-in
circuit has $O_\varphi(N^k)$ gates.

Encode every gate by its three-or-fewer Tseitin clauses and assert the
output. Each assignment to the relation-entry inputs determines exactly
one assignment to all gate outputs by evaluation in syntax order. It
satisfies the output assertion precisely when the corresponding
expansion satisfies $\varphi$. This proves parsimony, including the
unconstrained input variables for unused relation entries.

Name a gate by its subformula, assignment tuple, and its position in a
quantifier's binary bracketing. These fields use $O_\varphi(\log N)$
bits. All child indices are obtained by tuple restriction, replacement
of one coordinate, and counter arithmetic on fixed-length tuples.
Input-structure atoms are queried by scanning the supplied table for
the requested tuple if necessary. Time is not charged to the decoder
at this step. Thus every requested clause is accessible in the claimed
space.
\end{proof}

\begin{theorem}[Counting upper bound]
\label{thm:fo-counting-upper}
Every property in $\mathcal K_k$ can be decided simultaneously in
\[
 \mathrm{DSPACE}\left[O\!\left(\frac{N^k}{\sqrt{\log N}}\right)\right]
 \quad\hbox{and time}\quad
 2^{O(N^k/\sqrt{\log N})}.
\]
The constants may depend on the fixed formulas and vocabulary.
\end{theorem}
\begin{proof}
Apply the Boolean-network counting theorem of~\cite{Cruz26} to the
logspace-accessible Tseitin instance of Lemma~\ref{lem:fo-grounding}. If some relation inputs
are unused, include their unary weight vertices as in the weighted
network construction; alternatively factor their free choices into
an explicitly known power of two. The network still has $O(N^k)$
edges, so requested count bits use the stated bounds.

There are $O(N^k)$ guessed relation bits, hence every count has
$O(N^k)$ output bits. Compare two counts from their most significant
bits down, querying them afresh one bit at a time. Store the current
position in $O(\log N)$ bits, and stop at the first unequal pair.
For a fixed modulus, update a fixed-size residue accumulator while
scanning the bits. The polynomial number of count-bit calls does not
change the displayed exponential running-time rate. A fixed Boolean
combination costs a constant number of such computations.
\end{proof}

The lower bound side is a standard tally diagonalization at the scale
$N^k/\sqrt{\log N}$.

\begin{lemma}[Unary space separation at this scale]
\label{lem:tally-separation}
For every fixed $k\geq1$ there is a tally language $L_k$ in
$\mathrm{DSPACE}[N^k]$ but not in
$\mathrm{DSPACE}[N^k/\sqrt{\log N}]$.
\end{lemma}
\begin{proof}
Enumerate all deterministic machines with one read-only endmarker-bounded
input tape and one binary work tape by binary descriptions $M_i$.
Every fixed finite multitape machine can be encoded in this model with
constant-factor space overhead: use fixed-length blocks for symbols
and head markers and sweep to simulate a transition.

On input $1^N$, put $i=\nu_2(N)$, with a fixed default index when $i=0$.
Each fixed positive $i$ occurs for infinitely many unbounded values of
$N$. A description of $M_i$ can be recovered from its index using
$O(\log N)$ space. Set
\[
 b(N)=\left\lfloor
 \frac{N^k}{\lceil\log_2\log_2(N+4)\rceil}
 \right\rfloor.
\]
Simulate $M_i$ on $1^N$ with at most $b(N)$ work cells. Reject a
simulation that crosses this boundary, and stop a simulation after a
configuration-count bound. If $q_i$ is its state count, at most
$q_i(N+2)(b(N)+1)2^{b(N)}$ configurations are possible; the logarithm
of this bound is $b(N)+O(\log N)$, so its counter fits in $O(b(N)+\log N)$
space. Define $L_k$ to accept exactly when this bounded simulation does
not accept. Thus $L_k\in\mathrm{DSPACE}[N^k]$.

Suppose a fixed $M_i$ decides $L_k$ in at most
$cN^k/\sqrt{\log N}$ cells. Since
$b(N)/(N^k/\sqrt{\log N})\to\infty$, that machine fits in the
simulation bound for every sufficiently large $N$. Choose such an $N$
with $\nu_2(N)=i$. The machine halts, and it cannot repeat a
configuration before halting. Therefore the bounded simulation returns
its actual answer, while the definition of $L_k$ returns its opposite.
This is a contradiction.
\end{proof}

\paragraph{Vectorial fixed points.}
Let $\mathrm{vPFP}_k$ permit simultaneous partial fixed points of a
fixed finite number of relation variables, each of arity at most $k$,
and Boolean flags. Its first-order update formulas have no variable
bound. The iteration starts empty, yields its limiting tuple if a
fixed point is reached, and yields empty relations otherwise.

\begin{theorem}[Same-arity logical separation]
\label{thm:same-arity}
On ordered structures of maximum input-relation arity at most $k$,
\[
 \mathcal K_k\subsetneq\mathrm{vPFP}_k.
\]
The noncontainment is already witnessed by a property of the cardinality
of a pure finite order.
\end{theorem}
\begin{proof}
A fixed tuple of $k$-ary relations has $O(N^k)$ bits. Iteration of its
fixed first-order update formulas uses two copies of that tuple, a
configuration counter of $O(N^k)$ bits, and $O(\log N)$ workspace to
evaluate each fixed formula. Nested fixed points multiply the space
by a fixed constant. Thus $\mathrm{vPFP}_k$ uses $O(N^k)$ space.

Conversely, a deterministic machine using $cN^k$ work cells is encoded
by a fixed number of tracks indexed by $[N]^k$ in lexicographic order.
Unary tape symbols in the machine become $k$-ary predicates, head
positions become singleton $k$-ary predicates, and machine states become
flags. Lexicographic successor on $k$-tuples is first-order definable
from the domain order: a successor changes one coordinate to its
successor, resets the following coordinates to minimum, and leaves
preceding coordinates unchanged. There are $k$ cases in this fixed
definition. Input heads for relation-table encodings use the same tuple
representation with a finite flag for the input symbol's relation.
First-order conditions retrieve the corresponding input-relation bit.
The finite transition table therefore yields fixed first-order updates.
Initialize by a flag and make halting configurations stationary.
This defines every such deterministic space computation in
$\mathrm{vPFP}_k$.

Theorem~\ref{thm:fo-counting-upper} now gives the inclusion. Apply
Lemma~\ref{lem:tally-separation} to $L_k$, identifying $1^N$ with the
pure order on $[N]$. Its machine has a fixed $\mathrm{vPFP}_k$
definition. If this property belonged to $\mathcal K_k$,
Theorem~\ref{thm:fo-counting-upper} would contradict
Lemma~\ref{lem:tally-separation}. This proves strictness.
\end{proof}

\paragraph{Hand check.}
For $k=2$, take one guessed binary relation $R$ and
$\varphi=\forall x\exists y\,R(x,y)$. On a two-element domain each
row of $R$ has three permitted nonempty subsets, independently. Thus
$c_\varphi=3^2=9$. The grounding has four relation-input bits and
asserts $(R(1,1)\vee R(1,2))\wedge(R(2,1)\vee R(2,2))$; each of its
nine assignments has exactly one Tseitin extension.

\paragraph{Relation to the arity hierarchy problem.}
Durand, Grandjean and Olive~\cite[Introduction]{DGO04} explicitly
formulate the
open arity hierarchy question with \emph{no bound on the number of
first-order variables}. The bound above is not an upper bound for that
whole fragment. With $v>k$ first-order variables, a direct grounding may
have $\Theta(N^v)$ gates and clauses, giving
$O(N^v/\sqrt{\log N})$ space rather than $o(N^k)$.
Allowing quantified functions also changes the interface: a general
$k$-ary function has $N^k\log N$ bits, not $N^k$.

\paragraph{A stronger target.}
For the universal, quantifier-free relational fragment with guessed
arity at most $k$, grounding produces a bounded-width CNF on
$O(N^k)$ guessed bits but potentially $N^v$ clauses. A uniform
\emph{variable-parameter} $o(n)$-space algorithm for bounded-width CNF
counting, with polynomial-length read-only input, would remove this
particular first-order-variable obstruction for that fragment.
The known subexponential variable-parameter time/formula-size bounds
from sparsification do not by themselves establish this space claim:
the branch generator and its leaf-instance decoder must also fit in
$o(n)$ space. We state this as the missing lemma; it is not asserted
as a consequence of anything proved here.

\paragraph{Related work.}
The conference abstract of Gr\"adel and Rosen~\cite{GR99} states the
regular-language characterization for $\Sigma^1_1(\mathrm{FO}^2)$.
Grandjean~\cite[Theorem~0.1 and introduction]{Gra84} distinguishes
symbol arity from the number of universal quantifiers; function
symbols are permitted there, unlike in our relational grounding
lemma. Grandjean and Olive~\cite{GO04} use unit-cost nondeterministic
RAMs, so their capture theorem cannot simply be substituted for a
bit-cost multitape Turing-machine bound.


\section{A conditional route to subexponential Frege refutations}
\label{sec:few-order-frege-milestone}

This section isolates what is missing for Frege. The reflection
hypothesis~\eqref{eq:frege-reflection-hypothesis} below is
\emph{conditional}; we do not prove it. The size-transfer implications
from that hypothesis are proved. The hypothesis is a target for the
specific contraction certificates, not a claimed simulation of
arbitrary commutative IPS. Section~\ref{sec:frege-obstacles} explains
why the known simulations of IPS by Frege do not supply it.

\subsection{The reflection hypothesis and the size transfer}

\paragraph{Proof model and certificate promise.}
Fix one finite Frege system over the connectives $\neg,\wedge,\vee$,
with the usual substitution of formulas for schematic variables.
A refutation of a CNF $F$ means a Frege proof of the tautology
$\neg F$. Proof size is total written symbol size, counting all
formula copies; variable names may be written in binary.
Put $s=n+m+2$ for a 3-CNF with $n$ specified variables and $m$
clauses. Its bit encoding has size $O(s\log(s+2))$.

For every allowed power of two $h$ with $2^{3h}\le s$, the
constructions of Sections~\ref{sec:few-order} and~\ref{sec:ips-universal}
give an integral Hilbert-like IPS
certificate $C(x,y)$ of total formula size
\[
 S\le 2^{K s/\sqrt h}
 \tag{F1}\label{eq:frege-compiled-size}
\]
for an absolute constant $K$. It is a sum of at most $h$ ordered,
syntactically multilinear formulas. The orders belong to the
universal list arising from the global row/column histories and
prescribed owner placement. Its binary product depth is
$O(\log(s+2))$, formal degree is at most $n$, and placeholder
degree is at most one at every gate. One can obtain all these
properties simultaneously by compiling each first-failed-clause
coefficient with the same $h$ and owner groups, putting all
placeholders before the original variables, and grouping the
coefficients by history. The factor $m$ in size is absorbed in
\eqref{eq:frege-compiled-size}; it adds no product-depth factor.
Reduction modulo two gives a certificate over $\mathbb F_2$.
These explicitly generated certificates, together with their
construction data, form the promised input class below.

\paragraph{Conditional reflection hypothesis.}
There are constants $A,c,C>0$ such that every promised size-$S$
certificate with order parameter $h$ for an unsatisfiable 3-CNF
$F$ has a Frege refutation of size at most
\[
 S^{A(1+\log(h+1))^c}\,
       2^{A(\log(s+2))^C}.
 \tag{F2}\label{eq:frege-reflection-hypothesis}
\]
This hypothesis includes the necessary internal proof of soundness
or identity implications. An external identity-testing algorithm
with this running time alone is not asserted to imply
\eqref{eq:frege-reflection-hypothesis}. The hypothesis is only a
proof-size statement; no automatization claim is included.

\begin{proposition}[Conditional subexponential Frege transfer]
\label{prop:conditional-frege-transfer}
If \eqref{eq:frege-reflection-hypothesis} holds, every unsatisfiable
3-CNF with parameter $s=n+m+2$ has a Frege refutation of size
\[
 2^{\,O\!\left(
      \dfrac{s(\log\log(s+3))^c}{\sqrt{\log(s+2)}}
        +(\log(s+2))^C\right)}
 =2^{o(s)}.
\]
In particular, this gives $2^{o(n)}$ Frege refutations when
$m=O(n)$.
\end{proposition}

\begin{proof}
For all sufficiently large $s$, choose $h$ to be the largest power
of two at most $(\log_2s)/3$. Then $2^{3h}\le s$, and rounding
down to a power of two loses a factor smaller than two. Hence
$h=\Theta(\log s)$. Taking base-two logarithms in
\eqref{eq:frege-reflection-hypothesis} and using
\eqref{eq:frege-compiled-size} gives, for the resulting refutation
$\mathcal P_F$ of $F$,
\[
 \log_2|\mathcal P_F|
 \le AK\frac{s}{\sqrt h}(1+\log(h+1))^c
                  +A(\log(s+2))^C.
\]
The first term is
$O(s(\log\log s)^c/\sqrt{\log s})$.
To see that it is $o(s)$, put $u=\log s$. Its ratio to $s$ is
$O((\log u)^c/u^{1/2})$, which tends to zero: after taking
logarithms of that ratio, the negative term
$-\tfrac12\log u$ dominates $c\log\log u$.
The second term is $o(s)$ because every fixed power of $\log s$
is $o(s)$. This proves the claimed bound. If $m=O(n)$, then
$s=O(n)$ and $s\ge n$, so $o(s)=o(n)$.
\end{proof}

\begin{proposition}[The exponent-loss threshold]
\label{prop:frege-loss-threshold}
Suppose instead that the promised certificates admit Frege
refutations of size
\[
 S^{p(h)}2^{O((\log(s+2))^C)}
\]
for a nonnegative function $p$ and a fixed $C$.
For any allowed choice $h=h(s)\to\infty$, this bound and
\eqref{eq:frege-compiled-size} imply $2^{o(s)}$ refutations whenever
\[
 \frac{p(h(s))}{\sqrt{h(s)}}\longrightarrow0.
\]
In particular, $p(h)=o(\sqrt h)$ suffices. This is the threshold
for preserving the saving by substitution in these bounds; it is
not a lower bound ruling out a sharper simulation.
\end{proposition}

\begin{proof}
The logarithm of the resulting size bound is
\[
 O\!\left(s\frac{p(h(s))}{\sqrt{h(s)}}
                  +(\log(s+2))^C\right).
\]
Its ratio to $s$ tends to zero under the displayed assumption.
If $p(h)$ is only bounded by a constant multiple of $\sqrt h$,
substitution gives only an $O(s)$ exponent. A larger displayed loss
does not yield a vanishing ratio from this computation. This last
observation is a limitation of the available upper-bound estimate,
not a claim that the actual shortest Frege proof must be large.
\end{proof}

\subsection{Why known simulations do not supply the hypothesis}
\label{sec:frege-obstacles}

\paragraph{Identity testing for sums of ROABPs.}
Gurjar, Korwar, Saxena and Thierauf~\cite[Theorem~3.2 and the explicit
bound in the introduction]{GKST15} give white-box time
$(Nd\,w^{2^r})^{O(r)}$ for a sum of $r$ ROABPs on $N$ variables,
individual degree $d$, and width $w$.
With $r=h$ and parameters bounded polynomially by $S$, the displayed
time bound is $S^{O(h2^h)}$. The corresponding exponent-loss
estimate divided by $\sqrt h$ grows like $\sqrt h\,2^h$ instead of
tending to zero. Thus directly substituting that bound cannot
give the desired saving, even if a matching Frege formalization
were available. Choosing $h$ to grow more slowly does not repair
this displayed loss. Keeping $h$ constant leaves the original
certificate-size exponent proportional to $s$.

\paragraph{The two-bucket simulation.}
Chatterjee, Chatterjee, Ghosal and Mukhopadhyay~\cite[Definition~1.1
and Theorems~1.2 and~3.3]{CCGM26} simulate a class of IPS certificates
in Frege. Their partition of the variables into two buckets is global.
Variables in the same bucket do not commute; variables in distinct
buckets do commute. Placeholder variables receive no extra commutation
relations. The general Frege simulation is quasipolynomial in
certificate size. The following proposition shows that the history
orders do not fit a fixed bucket system with few buckets.

\begin{proposition}[Antidiagonal obstruction]
\label{prop:antidiagonal}
Consider a $g\times g$ array of distinct formal symbols $X_{a,b}$.
Suppose row-major and column-major products of all symbols are to be
identified solely through a global bucket system in which symbols
commute only across different buckets. Then at least $g$ buckets are
necessary.
\end{proposition}

\begin{proof}
The antidiagonal symbols $X_{1,g},X_{2,g-1},\ldots,X_{g,1}$ occur in
that order in the row-major word and in the reverse order in the
column-major word. In a partially commutative word algebra, projection
onto the symbols of one bucket is invariant under all permitted swaps:
swapping symbols from different buckets does not change either
within-bucket projected word. If two antidiagonal symbols belonged to
one bucket, its two projected words would put that pair in opposite
orders. They would therefore differ, so the original words would not
be equal. Each antidiagonal symbol needs a different bucket.
\end{proof}

For $g=3$, the symbols $X_{1,3},X_{2,2},X_{3,1}$ already require
three buckets. This proves failure of the literal two-bucket lifting.
It does not rule out every proof with additional commutator
certificates, and constructing small certificates for the required
reorderings remains a possible route.

\paragraph{Quasipolynomial dependence on certificate size.}
A bound $S^{(\log S)^b}$ with any fixed $b>0$ substitutes to
\[
 2^{O((s/\sqrt h)^{b+1})}.
\]
Even the maximal allowed $h=\Theta(\log s)$ gives an exponent
estimate whose ratio to $s$ is
$s^b/(\log s)^{(b+1)/2}$, which diverges. In the same way, if a
simulation costs $S^{(\log S)^c}$ and $S=2^{O(n/\sqrt{\log n})}$,
substitution alone yields $2^{O((n/\sqrt{\log n})^{c+1})}$, which is
not subexponential in $n$ for any fixed $c\ge1$. The two-bucket
simulation~\cite[Theorems~1.2 and~3.3]{CCGM26} has such a
quasipolynomial dependence on certificate size. It also has further
quasipolynomial ABP-to-formula and witness costs after its initial
homogenization~\cite[Appendices~B.4 and~B.6]{CCGM26}. Thus observing
that the original weight polynomial is homogeneous is not enough to
make that theorem preserve the counting saving. Its fixed global
two-bucket commutation model also does not directly match the history
orders, by Proposition~\ref{prop:antidiagonal}. Citing that theorem
therefore supplies neither the structural applicability nor the
needed quantitative reflection.

\paragraph{A useful positive refinement.}
The path-index lemma of~\cite{Cruz26} bounds the sum of product-index
lengths along a path by $O(\log s)$. After binary expansion, the
product depth is therefore $O(\log s)$ even though the formula size is
exponential in $s/\sqrt{\log s}$. Degree-tracking homogenization of
such a formula has at most $(n+1)^{O(\log s)}$ copies per node. For
$n\le s$ this is $2^{O(\log^2s)}$, absorbed by the original
formula-size rate. This removes one potential homogenization loss. It
does not supply the missing global noncommutative lift or remove the
later witness costs in the two-bucket Frege simulation.

\subsection{Scope and a numerical check}

\paragraph{Scope of the variable bound.}
The conclusion above is $2^{o(n+m)}$, and becomes $2^{o(n)}$ for
linear-size CNFs. Counting sparsification alone does not justify
extending this Frege conclusion to arbitrary 3-CNFs in the number
of variables. Such an extension would additionally need a Frege
derivation that combines the sparse-branch refutations with a
proof of the sparsification decomposition, within the required
size. No such proof-efficient sparsification step is assumed or
proved here.

\paragraph{A numerical check.}
For $s=2^{12}$ the largest admissible power-of-two name length is
$h=4$, with two binary cut decisions and four history orders.
Substitution changes the certificate logarithmic-size bound
$Ks/2$ to at most
$AKs(1+\log5)^c/2+A(\log(s+2))^C$. This verifies the
parameter substitution on one actual allowed recursion size;
the asymptotic saving follows from the proved limit, not from
this finite check.

\section{Open problems}
\label{sec:open-map}

\paragraph{1. Frege reflection for the actual history orders.}
We regard this as the most important next target. The conditional
implication of Section~\ref{sec:few-order-frege-milestone} needs a loss
controlled by the number of merge levels, rather than by the logarithm
of the full expanded certificate size. The sufficient bound is
$S^{\operatorname{polylog}h}2^{\operatorname{polylog}s}$.
A first step would be a proof of the row/column switching identity
with explicit certificates for reordering its variable factors; a
second would be a composition theorem across all global histories that
does not distribute the full assignment sum. There is also an explicit
criterion for when to stop: by Proposition~\ref{prop:frege-loss-threshold},
a construction whose proved exponent loss is $\Omega(\sqrt h)$ does
not yield a subexponential bound through the transfer established
here. This is not an assertion that all such proofs must be large.

The October 2026 two-bucket Frege paper~\cite{CCGM26} is directly
relevant but does not establish the needed lift. Its commutation
partition is global, and its size dependence is quasipolynomial in
$S$. Proposition~\ref{prop:antidiagonal} and the substitution
calculation of Section~\ref{sec:frege-obstacles} identify two distinct
failures of applying it verbatim.

A natural approach is a proof of the contraction identities in Frege
that tracks the global merge history. By the results of
Section~\ref{sec:few-order}, only a slowly growing family of highly
structured orders is needed, so the proof no longer has to organize
the full assignment space. A concrete first goal would be a
reordering certificate for one merge that composes with loss
polynomial in the number of merge levels. A bound polynomial in
$\log h$, combined with the conditional transfer of
Proposition~\ref{prop:conditional-frege-transfer}, would turn the
counting construction into subexponential Frege proofs for every
linear-size 3-CNF.

\paragraph{2. The quantitative size--order tradeoff.}
The qualitative threshold is proved, but the number of orders needed
at size $2^{O(N/\sqrt{\log N})}$ is only known to lie between
$\Omega(\log\log N)$ and $O(\log N)$
(Theorem~\ref{thm:fast-code-order-separation}). Closing this interval
is a focused open question. More generally, the exponents
$\Omega(N/2^k)$ and $O(N/\sqrt k)$ are far apart.
The Homeister-style proof loses a factor two in residual block size
each time it accommodates another order. The compiler loses a square
root when potential is converted to index length. Neither calculation
identifies the true tradeoff. A first step would be to analyze the
specific recursively permuted order list instead of arbitrary
permutations. No sharper tradeoff is claimed here.

\paragraph{3. A few-order lower bound strong enough for Boolean circuits.}
Proposition~\ref{prop:coefficient-bridge} gives a precise target:
$N^{\Omega(d)}$ lower bounds against logarithmically many
different-order set-multilinear ABPs for a Boolean coefficient
polynomial. This would force
$\Omega(b\sqrt{\log b})$ Boolean circuit size on its
$b=d\log N$-bit coefficient predicate. The code lower bound of
Section~\ref{sec:parity-findings} cannot do this for linear-size
functions, since the compiler itself supplies small enough
representations once the number of orders grows.
The bound of~\cite{CKSS24} has the required scale only
for a single order; its sum-of-orders theorem is quantitatively weaker.
The 2026 min-partition-rank paper~\cite{FLSY26} is relevant to the
choice of measure; its full-rank examples prevent treating full rank
alone as a sufficient lower-bound argument.

\paragraph{4. Extending the program-proof lower bound.}
Theorem~\ref{thm:ka-tree-superlinear} measures total symbols in a
fully written proof tree. The next tasks are to establish a comparable
verification rate for an ordinary shared proof representation, and to
construct an explicit family of valid identities with a quantitative
proof-length lower bound. A verifier costing $O(L\log L)$ already
uses more than the available logarithmic saving when $L\ge n$, so
the present transfer does not automatically extend to
arbitrary binary line references. Existing linear-checkable
certificates, if any, should be analyzed by their actual encoding.

The chapter of Beier and Holzer~\cite{BH24} concerns the proof
complexity of Salomaa's $F_1$. Its full text was not available to us;
its publisher metadata and a reproduced abstract were. Schweigler's
2025 thesis~\cite[Section~5, pp.~29--31]{Sch25} records general
polynomial-proof characterizations, not a quantitative tree-like
bound of the form proved here. This search does not establish a
complete priority claim for the Kleene-algebra result.

\paragraph{5. The first-order-variable obstruction in arity problems.}
The same-arity separation of Section~\ref{sec:bounded-variable}
requires an $\mathrm{FO}^k$ matrix. The classic arity question permits
arbitrarily many first-order variables; Durand, Grandjean and
Olive~\cite[Introduction, pp.~1--2]{DGO04} explicitly distinguish
those parameters. Guessing relations of arity $k$ uses $O(N^k)$ bits,
but grounding a $v$-variable matrix may require $N^v$ constraints.
When $v>k$, a clause-sensitive saving in $N^v$ does not become a
space saving below $N^k$. This is a polynomial gap in the parameters,
not a missing constant. A variable-parameter, uniformly sublinear-space
bounded-width counting theorem with polynomial-size read-only input
would address that interface; the existing nonuniform
sparsification argument does not supply it. We claim no solution to
Fagin's full arity hierarchy problem or to Asser's spectrum-complement
problem.

\paragraph{Related current work and attributions.}
The public list of STOC 2027 submissions~\cite{STOC27}, consulted on
19 November 2026, contains \#43, \emph{Hard Instances for Strong
Semantic Proof Systems: CNFs for Depth-Four PC and Deep Instances for
$\mathsf{AC}^0[2]$-Frege}; \#1046, \emph{Progress on AC0[2]-Frege Lower
Bounds}; \#1284, \emph{Fine-Grained Proof Complexity for Elementary
Arithmetic in Extended Frege}; and \#1702, \emph{The Contraction Lemma
Beyond ETH: Algorithms and Lower Bounds}. These titles show current
activity; they do not certify the contents or settle the novelty of
this paper. We checked the two-bucket Frege paper, the fast-code
theorems, the sum-of-ROABPs PIT bound, and the cited ABP lower bounds
in their primary texts. The code lower-bound mechanism is due to
Homeister~\cite{Homeister03}, the invalid-computation reduction is
classical~\cite{GS05}, and the first-failed-clause IPS identity follows
Grochow and Pitassi~\cite{GP18}.

\begin{thebibliography}{99}

\bibitem{Ant26} Claude (Anthropic).
\newblock A $2^{o(n)}$ time algorithm for 3-SAT.
\newblock Manuscript, 2026.
\newblock \url{https://drive.google.com/file/d/1u1XufQIU-PZafNWRO1Is-SUlX760agVp/view}.

\bibitem{BH24} Simon Beier and Markus Holzer.
\newblock On regular expression proof complexity of Salomaa's axiom
system~$F_1$.
\newblock In \emph{Taming the Infinities of Concurrency}, pages 72--100.
Springer, 2024.
\newblock \url{https://doi.org/10.1007/978-3-031-56222-8_4}.

\bibitem{BIR26} Martijn Brehm, Yuval Ishai, and Nicolas Resch.
\newblock Fast bounded-independence functions and their duals.
\newblock In \emph{ITC 2026}, LIPIcs 385, Article~8, 2026.
\newblock \url{https://doi.org/10.4230/LIPIcs.ITC.2026.8}.

\bibitem{BR26} Martijn Brehm and Nicolas Resch.
\newblock Linear time encodable binary code achieving GV bound with
linear time encodable dual achieving GV bound.
\newblock In \emph{ITCS 2026}, LIPIcs 362, Article~28, 2026.
\newblock \url{https://doi.org/10.4230/LIPIcs.ITCS.2026.28}.

\bibitem{CCGM26} Abhranil Chatterjee, Prerona Chatterjee,
Utsab Ghosal, and Partha Mukhopadhyay.
\newblock Quasi-polynomial Frege simulation of IPS beyond
noncommutativity.
\newblock ECCC TR26-166, revision~1, 30 October 2026.
\newblock \url{https://eccc.weizmann.ac.il/report/2026/166/revision/1/download/}.

\bibitem{CKSS24} Prerona Chatterjee, Deepanshu Kush, Shubhangi Saraf,
and Amir Shpilka.
\newblock Lower bounds for set-multilinear branching programs.
\newblock In \emph{CCC 2024}. Corrected version: ECCC TR23-212,
revision~2, 17 February 2024.
\newblock \url{https://eccc.weizmann.ac.il/report/2023/212/revision/2/download/}.

\bibitem{Cruz26} Charlie Cruz.
\newblock Separating nondeterministic time from deterministic space.
\newblock Manuscript supplied by the author; source file
\texttt{separating-ntime-from-space.tex}, consulted 19 November 2026.

\bibitem{DGO04} Arnaud Durand, Etienne Grandjean, and Fr\'ed\'eric Olive.
\newblock New results on arity vs.\ number of variables.
\newblock Manuscript, 31 March 2004.
\newblock \url{https://webusers.imj-prg.fr/~arnaud.durand/papers/dgo.pdf}.

\bibitem{FLSY26} Th\'eo Bor\'em Fabris, Nutan Limaye,
Srikanth Srinivasan, and Amir Yehudayoff.
\newblock Multilinear algebraic branching programs and the
min-partition rank method.
\newblock In \emph{CCC 2026}, LIPIcs 383, Article~22, 2026.
\newblock \url{https://doi.org/10.4230/LIPIcs.CCC.2026.22}.

\bibitem{GR99} Erich Gr\"adel and Eric Rosen.
\newblock Two-variable descriptions of regularity.
\newblock In \emph{LICS 1999}, pages 14--23, 1999.
\newblock \url{https://lics.siglog.org/archive/1999/GradelRosen-TwoVariableDescript.html}.
\newblock (The word-structure result uses successor, endpoints and
letters; we do not extend it to a supplied full-order predicate.)

\bibitem{GS05} Gregor Gramlich and Georg Schnitger.
\newblock Minimizing NFA's and regular expressions.
\newblock In \emph{STACS 2005}, LNCS 3404, pages 399--411, 2005.
\newblock Extended version:
\url{https://citeseerx.ist.psu.edu/document?doi=cb9486498319bae7bf0d28e5960115e2571bf2f2&repid=rep1&type=pdf}.

\bibitem{Gra84} Etienne Grandjean.
\newblock The spectra of first-order sentences and computational
complexity.
\newblock \emph{SIAM J. Comput.}, 13(2):356--373, 1984.

\bibitem{GO04} Etienne Grandjean and Fr\'ed\'eric Olive.
\newblock Graph properties checkable in linear time in the number of
vertices.
\newblock \emph{J. Comput. System Sci.}, 68:546--597, 2004.

\bibitem{GP18} Joshua A. Grochow and Toniann Pitassi.
\newblock Circuit complexity, proof complexity, and polynomial identity
testing: The ideal proof system.
\newblock \emph{J. ACM}, 65(6), Article~37, 2018.
\newblock \url{https://doi.org/10.1145/3230742}.

\bibitem{GKST15} Rohit Gurjar, Arpita Korwar, Nitin Saxena, and
Thomas Thierauf.
\newblock Deterministic identity testing for sum of read-once
oblivious arithmetic branching programs.
\newblock In \emph{CCC 2015}, 2015.
\newblock \url{https://www.cse.iitb.ac.in/~rgurjar/papers/sumROABP.pdf}.

\bibitem{Homeister03} Matthias Homeister.
\newblock Lower bounds for the sum of graph-driven read-once parity
branching programs.
\newblock ECCC TR03-068, 2003.
\newblock \url{https://eccc.weizmann.ac.il/report/2003/068/download/}.

\bibitem{Kozen94} Dexter Kozen.
\newblock A completeness theorem for Kleene algebras and the algebra of
regular events.
\newblock \emph{Information and Computation}, 110(2):366--390, 1994.
\newblock Author version: \url{https://www.cs.cornell.edu/~kozen/Papers/ka.pdf}.

\bibitem{PreviousNote}
\newblock Logical consequences of integral counting.
\newblock Preceding note in this series, 19 November 2026; source file
\texttt{logic\_implications\_of\_integral\_counting.tex}.

\bibitem{Sch25} Manuel Schweigler.
\newblock Proof systems for regular expressions.
\newblock Bachelor's thesis, TU Wien, 5 March 2025.
\newblock \url{https://www.dmg.tuwien.ac.at/hetzl/teaching/b_schweigler.pdf}.

\bibitem{STOC27} STOC 2027 submissions.
\newblock Public list, consulted 19 November 2026. Submission numbers
denote listed submissions, not accepted papers.
\newblock \url{https://sigact-stoc.github.io/stoc2027/}.

\bibitem{Szeider26} Stefan Szeider.
\newblock Consequences of the ETH refutations for proof complexity and
knowledge compilation.
\newblock Version~2, 16 November 2026.
\newblock \url{https://doi.org/10.5281/zenodo.24299623}.
\newblock See Section~2, assumption~(A1), and Section~7, Open Problem~5.

\end{thebibliography}
\end{document}
